\SourcePage{111}
\SGAsection{4}{对应的配对；Lefschetz 公式}
\textbf{4.1。} 对 $i=1,2$，设 $X_i$ 是 $S$ 上的\Term{scheme}{概形}，$X=X_1\times_S X_2$，$p_i:X\to X_i$ 为\Term{projection}{投影}，$L_i\in\operatorname{ob}D_{ctf}(X_i)$，$K_{X_i}$ 是\Term{dualizing-complex}{对偶化复形}（1.3）。由 2.1.1 与\Term{evaluation-map}{求值箭头} 2.2.2 复合得到的\Term{canonical}{典范}\Term{arrow}{箭头}
\[p_i^*R\mathcal Hom(L_i,K_{X_i})\mathbin{\overset L\otimes}p_i^*L_i\longrightarrow p_i^*K_{X_i}\tag{4.1.1}\]
在 $X$ 上取\Term{tensor-product}{张量积}，再与 Künneth \Term{isomorphism}{同构} 1.7.6 复合，给出\Term{pairing}{配对}
\[(DL_1\mathbin{\overset L\otimes}_S L_2)\mathbin{\overset L\otimes}(L_1\mathbin{\overset L\otimes}_S DL_2)\longrightarrow K_X.\tag{4.1.2}\]
借助 3.1.1、3.1.2，可将其写为
\[R\mathcal Hom_S(L_1,L_2)\mathbin{\overset L\otimes}R\mathcal Hom_S(L_2,L_1)\longrightarrow K_X.\tag{4.1.3}\]
由此得到称为\Term{cup-product}{杯积}的\Term{pairing}{配对}
\[\langle\ ,\ \rangle:\operatorname{Hom}_S(L_1,L_2)\otimes\operatorname{Hom}_S(L_2,L_1)\longrightarrow H^0(X,K_X).\tag{4.1.4}\]
也可将 4.1.2 看作\Term{evaluation-map}{求值箭头} $DL_i\mathbin{\overset L\otimes}L_i\to K_{X_i}$ 的\Term{external-tensor-product}{外张量积}。结合 Künneth \Term{isomorphism}{同构} 2.2.4，可见 4.1.2 是“\Term{perfect}{完美}\Term{duality}{对偶}”，即通过\Term{adjunction-isomorphism}{伴随同构} 2.2.1，将\Term{complex}{复形} $DL_1\mathbin{\overset L\otimes}_S L_2$、$L_1\mathbin{\overset L\otimes}_S DL_2$ 中每一个等同于另一个的\Term{duality}{对偶}（取值于 $K_X$）。
当 $X_1=X_2=S$ 时，4.1.3 是任意\Term{ringed-topos}{环化拓扑斯}上的\Term{perfect-complex}{完美复形}的\Term{pairing}{配对}（SGA 6 I 7.8）的特例；依照（SGA 6 I 8.3）的记号，有
\[\langle u,v\rangle=\operatorname{Tr}(vu)=\operatorname{Tr}(uv).\tag{4.1.5}\]
\textbf{4.2。} 设 $c:C\to X$、$d:D\to X$，$e=c\times_Xd:E=C\times_XD\to X$。对 $P,Q\in\operatorname{ob}D_{ctf}(X)$，有\Term{canonical}{典范}\Term{arrow}{箭头}（通常不是\Term{isomorphism}{同构}）
\SourcePage{112}
\[c^!P\mathbin{\overset L\otimes}_X d^!Q\longrightarrow e^!(P\mathbin{\overset L\otimes}Q),\tag{4.2.1}\]
其定义方式与 1.7.3 相同，是复合\Term{arrow}{箭头}
\[e_!(c^!P\mathbin{\overset L\otimes}_Xd^!Q)\xrightarrow[\sim]{(1)}c_!c^!P\mathbin{\overset L\otimes}d_!d^!Q\xrightarrow{(2)}P\mathbin{\overset L\otimes}Q\]
的\Term{adjunction}{伴随}；这里 (1) 是 Künneth \Term{isomorphism}{同构}（SGA 4 XVII 5.4.2.2），(2) 是\Term{adjoint}{伴随箭头}的\Term{tensor-product}{张量积}。对 4.2.1 应用 $e_!$，再与刚引用的 Künneth \Term{isomorphism}{同构}复合，得到\Term{canonical}{典范}\Term{arrow}{箭头}
\[c_!c^!P\mathbin{\overset L\otimes}d_!d^!Q\longrightarrow e_!e^!(P\mathbin{\overset L\otimes}Q).\tag{4.2.2}\]
同样，将 Künneth \Term{arrow}{箭头} $c_*c^!P\mathbin{\overset L\otimes}d_*d^!Q\to e_*(c^!P\mathbin{\overset L\otimes}_Xd^!Q)$ 与对 4.2.1 应用 $e_*$ 得到的\Term{arrow}{箭头}复合，得到\Term{canonical}{典范}\Term{arrow}{箭头}
\[c_*c^!P\mathbin{\overset L\otimes}d_*d^!Q\longrightarrow e_*e^!(P\mathbin{\overset L\otimes}Q).\tag{4.2.3}\]
取 $P=DL_1\mathbin{\overset L\otimes}_SL_2$、$Q=L_1\mathbin{\overset L\otimes}_SDL_2$，将 4.2.2、4.2.3 与 4.1.2 及\Term{isomorphism}{同构} $e^!K_X=K_E$ 复合，并作 3.1.1、3.2.1 的等同，得到\Term{pairing}{配对}
\[c_!R\mathcal Hom(c_1^*L_1,c_2^!L_2)\mathbin{\overset L\otimes}d_!R\mathcal Hom(d_2^*L_2,d_1^!L_1)\longrightarrow e_!K_E,\tag{4.2.4}\]
\[c_*R\mathcal Hom(c_1^*L_1,c_2^!L_2)\mathbin{\overset L\otimes}d_*R\mathcal Hom(d_2^*L_2,d_1^!L_1)\longrightarrow e_*K_E,\tag{4.2.4$'$}\]
\[\langle\ ,\ \rangle:\operatorname{Hom}(c_1^*L_1,c_2^!L_2)\otimes\operatorname{Hom}(d_2^*L_2,d_1^!L_1)\longrightarrow H^0(E,K_E).\tag{4.2.5}\]
这些\Term{pairing}{配对}与\Term{etale-localization}{平展局部化}相容：给定交换方形（不要求为\Term{cartesian}{笛卡尔}方形）
\[\begin{tikzcd}[column sep=small] &E'\arrow[dl]\arrow[dr]&&&E\arrow[dl]\arrow[dr]&\\ C'\arrow[dr]&&D'\arrow[dl]&C\arrow[dr]&&D\arrow[dl]\\ &X'&&&X&\end{tikzcd}\]
左图\Term{etale}{平展}地位于右图之上。若以 $c',d',e'$ 分别记复合 $C'\to C\to X$、$D'\to D\to X$、$E'\to E\to X$，则有交换方形
\SourcePage{113}
\[\begin{tikzcd}[column sep=small]
\operatorname{Hom}(c_1^*L_1,c_2^!L_2)\otimes\operatorname{Hom}(d_2^*L_2,d_1^!L_1)\arrow[r,"4.2.5"]\arrow[d]&H^0(E,K_E)\arrow[d]\\
\operatorname{Hom}(c_1'^*L_1,c_2'^!L_2)\otimes\operatorname{Hom}(d_2'^*L_2,d_1'^!L_1)\arrow[r]&H^0(E',K_{E'})
\end{tikzcd}\tag{4.2.6}\]
竖直\Term{arrow}{箭头}是\Term{restriction}{限制}，下方水平\Term{arrow}{箭头}是类似于 4.2.5 的\Term{arrow}{箭头}与\Term{restriction}{限制} $H^0(E'',K_{E''})\to H^0(E',K_{E'})$ 的复合，其中 $E''=C'\times_{X'}D'$（对 4.2.4、4.2.4$'$ 型\Term{arrow}{箭头}也有类似交换方形，留给读者写出）。若 $Z$ 在 $E$ 上\Term{etale}{平展}，则对 $u\in\operatorname{Hom}(c_1^*L_1,c_2^!L_2)$、$v\in\operatorname{Hom}(d_2^*L_2,d_1^!L_1)$，以
\[\langle u,v\rangle_Z\in H^0(Z,K_Z)\tag{4.2.7}\]
记 $\langle u,v\rangle$ 经\Term{restriction}{限制}在 $H^0(Z,K_Z)$ 中的像（实际应用中，$Z$ 将是 $E$ 的一个\Term{connected-component}{连通分支}）。若 $z$ 是 $E$ 的孤立\Term{closed-point}{闭点}，$\bar z$ 是其上的\Term{geometric-point}{几何点}，则由 4.2.6 可知，$\langle u,v\rangle_{\bar z}\in H^0(\bar z,K_{\bar z})=A$ 只依赖于在 $\bar z$ 于 $C,D,X$ 中的像处作\Term{strict-localization}{严格局部化}所得的情形。
\begin{SGAresult}{例}{4.3}
假设 $X_i$ \Term{smooth}{光滑}，纯\Term{relative-dimension}{相对维数}为 $r_i$，$c,d$ 是\Term{closed-immersion}{闭浸入}，而 $c_2$（相应地 $d_1$）\Term{finite}{有限}。对 $L_1=A_{X_1}$、$L_2=A_{X_2}$，由 3.1.3 有
\[DL_1\mathbin{\overset L\otimes}_SL_2=A_X(r_1)[2r_1],\qquad L_1\mathbin{\overset L\otimes}_SDL_2=A_X(r_2)[2r_2],\]
而 4.1.2 等同于乘积
\[A_X(r_1)[2r_1]\mathbin{\overset L\otimes}A_X(r_2)[2r_2]\longrightarrow A_X(r)[2r],\qquad r=r_1+r_2.\]
因此结合 3.2.1，4.2.5 等同于通常的\Term{cup-product}{杯积}
\[H_C^{2r_1}(X,A(r_1))\otimes H_D^{2r_2}(X,A(r_2))\longrightarrow H_{C\cap D}^{2r}(X,A(r)).\]
以 $c$ 的类（相应地 $d$ 的类）（3.2 a)）作为 $X_1$ 到 $X_2$（相应地 $X_2$ 到 $X_1$）的\Term{cohomological-correspondence}{上同调对应}。由（SGA $4\frac12$ Cycle 2.3.8）知，在 $C\cap D$ 的\Term{isolated-point}{孤立点} $z$ 处，
\SourcePage{114}
$\langle\operatorname{cl}(c),\operatorname{cl}(d)\rangle_z$ 正是 $C$ 与 $D$ 在 $z$ 处的\Term{intersection-multiplicity}{交数}在 $A$ 中的像。
\end{SGAresult}
下面给出本讲的主要结果。
\begin{SGAresult}{\Term{theorem}{定理}}{4.4}
对 $i=1,2$，设 $f_i:X_i\to Y_i$ 为\Term{proper}{固有} $S$ \Term{morphism}{态射}，$f=f_1\times_Sf_2:X\to Y$，$p_i:X\to X_i$、$q_i:Y\to Y_i$ 为\Term{projection}{投影}。再设有水平面为\Term{cartesian}{笛卡尔}方形的交换立方体
\[\begin{tikzcd}[row sep=large,column sep=large]
C'\arrow[dd,"f'"']\arrow[dr,"c'"]&&C\arrow[ll]\arrow[dd,"g"]\arrow[dr]\arrow[dl,"c" near end]\\
&X\arrow[dd,"f" near start]&&C''\arrow[ll,"c''"']\arrow[dd,"f''"]\\
D'\arrow[dr,"d'"']&&D\arrow[ll]\arrow[dr]\arrow[dl,"d" near end]\\
&Y&&D''\arrow[ll,"d''"]
\end{tikzcd}\tag{4.4.0}\]
其中 $c',c'',d',d''$ \Term{proper}{固有}。置 $c_i'=p_ic'$、$c_i''=p_ic''$、$d_i'=q_id'$、$d_i''=q_id''$。最后设 $L_i\in\operatorname{ob}D_{ctf}(X_i)$。则有交换方形
\[\begin{tikzcd}[column sep=small]
f_*c'_*R\mathcal Hom(c_1'^*L_1,c_2'^!L_2)\mathbin{\overset L\otimes}f_*c''_*R\mathcal Hom(c_2''^*L_2,c_1''^!L_1)\arrow[r,"(1)"]\arrow[d,"(2)"']&f_*c_*K_C\arrow[d,"(4)"]\\
d'_*R\mathcal Hom(d_1'^*M_1,d_2'^!M_2)\mathbin{\overset L\otimes}d''_*R\mathcal Hom(d_2''^*M_2,d_1''^!M_1)\arrow[r,"(3)"]&d_*K_D
\end{tikzcd}\tag{4.4.1}\]
其中 $M_i=f_{i*}L_i$；(3) 由 4.2.4 给出，(1) 由对 4.2.4 应用 $f_*$ 并与\Term{canonical}{典范}\Term{arrow}{箭头} $f_*c'_*\mathbin{\overset L\otimes}f_*c''_*\to f_*(c'_*\mathbin{\overset L\otimes}c''_*)$ 复合得到；(2) 是与含 $f$ 的两个面对应的 3.7.5 \Term{arrow}{箭头}的\Term{tensor-product}{张量积}；(4) 由\Term{adjoint}{伴随箭头} $g_*K_C=g_*g^!K_D\to K_D$ 定义。
\end{SGAresult}
\begin{proof}
置 $P=DL_1\mathbin{\overset L\otimes}_SL_2$、$Q=L_1\mathbin{\overset L\otimes}_SDL_2$、$E=DM_1\mathbin{\overset L\otimes}_SM_2$、$F=M_1\mathbin{\overset L\otimes}_SDM_2$。考虑图
\SourcePage{115}
\[\begin{tikzcd}[column sep=small]
f_*c'_*c'^!P\mathbin{\overset L\otimes}f_*c''_*c''^!Q\arrow[rr,"(1)"]\arrow[d,"(8)"']&&f_*c_*c^!(P\mathbin{\overset L\otimes}Q)\arrow[r,"(2)"]\arrow[d,"(9)"]&f_*c_*c^!K_X\arrow[d,"(10)"]\\
d'_*d'^!f_*P\mathbin{\overset L\otimes}d''_*d''^!f_*Q\arrow[r,"(3)"]\arrow[d,"(11)"']&d_*d^!(f_*P\mathbin{\overset L\otimes}f_*Q)\arrow[r,"(4)"]\arrow[d,"(12)"]&d_*d^!f_*(P\mathbin{\overset L\otimes}Q)\arrow[r,"(5)"]&d_*d^!f_*K_X\arrow[d,"(13)"]\\
d'_*d'^!E\mathbin{\overset L\otimes}d''_*d''^!F\arrow[r,"(6)"]&d_*d^!(E\mathbin{\overset L\otimes}F)\arrow[rr,"(7)"]&&d_*d^!K_Y
\end{tikzcd}\tag{4.4.2}\]
各\Term{arrow}{箭头}如下：(1) 由对 4.2.2 应用 $f_*$ 得到；(2) 由 4.1.2 定义；(3) 由 4.2.2 给出；(4) 由\Term{canonical}{典范}\Term{arrow}{箭头} $f_*P\mathbin{\overset L\otimes}f_*Q\to f_*(P\mathbin{\overset L\otimes}Q)$ 给出；(5) 由 4.1.2 给出；(6) 由 4.2.2 给出；(7) 由 4.1.2 给出；(8) 为 3.7.3 \Term{arrow}{箭头}的\Term{tensor-product}{张量积}；(9)、(10) 由 3.7.3 定义；(11)、(12) 由\Term{isomorphism}{同构} $f_*P\to E$、$f_*Q\to F$（3.3.1）给出；(13) 由\Term{adjoint}{伴随箭头} $f_*K_X\to K_Y$ 给出。由定义，经 3.1.1、3.2.1，4.4.1 等同于 4.4.2 的外\Term{boundary}{边界}。因此只须证明方形 (A)、(B)、(C)、(D) 交换，其中 (A) 为上排左部，(B) 为上排右部，(C) 为下排左部，(D) 为下排右部。(B)、(C) 的\Term{commutativity}{交换性}由\Term{functoriality}{函子性}显然得到；余下证明 (A)、(D)。

 a) (D) 的\Term{commutativity}{交换性}。只须证明对应于 $d=\operatorname{Id}$ 的方形
\[\begin{tikzcd}
f_*P\mathbin{\overset L\otimes}f_*Q\arrow[r]\arrow[d]&f_*K_X\arrow[d]\\ E\mathbin{\overset L\otimes}F\arrow[r]&K_Y
\end{tikzcd}\tag{4.4.3}\]
交换。回忆\Term{pairing}{配对} $P\mathbin{\overset L\otimes}Q\to K_X$（相应地 $E\mathbin{\overset L\otimes}F\to K_Y$）是\Term{evaluation-map}{求值箭头} $DL_i\mathbin{\overset L\otimes}L_i\to K_{X_i}$（相应地 $DM_i\mathbin{\overset L\otimes}M_i\to K_{Y_i}$）的\Term{external-tensor-product}{外张量积}。经 Künneth \Term{isomorphism}{同构} $f_{1*}DL_1\mathbin{\overset L\otimes}_Sf_{2*}L_2\xrightarrow\sim f_*P$、$f_{1*}L_1\mathbin{\overset L\otimes}_Sf_{2*}DL_2\xrightarrow\sim f_*Q$，4.4.3 等同于以下方形的\Term{external-tensor-product}{外张量积}：
\[\begin{tikzcd}
f_{i*}DL_i\mathbin{\overset L\otimes}f_{i*}L_i\arrow[r,"(1)"]\arrow[d,"(3)"']&f_{i*}K_{X_i}\arrow[d,"(4)"]\\ DM_i\mathbin{\overset L\otimes}M_i\arrow[r,"(2)"]&K_{Y_i}.
\end{tikzcd}\tag{4.4.4}\]
这里 (2) 是\Term{evaluation-map}{求值箭头}，(1) 是\Term{canonical}{典范}\Term{arrow}{箭头}
\SourcePage{116}
$f_{i*}DL_i\mathbin{\overset L\otimes}f_{i*}L_i\to f_{i*}(DL_i\mathbin{\overset L\otimes}L_i)$ 与求值所定义\Term{arrow}{箭头}的复合，(3) 是\Term{duality-isomorphism}{对偶同构} $f_{i*}DL_i\to DM_i$ 与 $M_i$ 的恒等\Term{arrow}{箭头}的\Term{tensor-product}{张量积}，(4) 是\Term{adjoint}{伴随箭头}。因此只须核对 4.4.4 交换。上述\Term{duality-isomorphism}{对偶同构}是\Term{canonical}{典范}\Term{arrow}{箭头}（SGA 4 XVIII 3.1.9.5）
\[f_{i*}R\mathcal Hom(L_i,K_{X_i})\longrightarrow R\mathcal Hom(f_{i*}L_i,f_{i*}K_{X_i})\]
与由\Term{adjoint}{伴随箭头} $f_{i*}K_{X_i}\to K_{Y_i}$ 得到的\Term{arrow}{箭头}的复合，故只须证明
\[\begin{tikzcd}
f_{i*}R\mathcal Hom(L_i,K_{X_i})\mathbin{\overset L\otimes}f_{i*}L_i\arrow[r,"(1)"]\arrow[d,"(3)"']&f_{i*}K_{X_i}\arrow[d,equal]\\ R\mathcal Hom(f_{i*}L_i,f_{i*}K_{X_i})\mathbin{\overset L\otimes}f_{i*}L_i\arrow[r,"(2)"]&f_{i*}K_{X_i}
\end{tikzcd}\tag{4.4.5}\]
交换，其中 (1) 为 4.4.4 的 (1)，(2) 为求值，(3) 为刚才的\Term{canonical}{典范}\Term{arrow}{箭头}与 $f_{i*}L_i$ 的\Term{tensor-product}{张量积}。更一般地，对 $D_{ctf}(X_i)$ 中任意\Term{arrow}{箭头} $N_i\to f_{i*}L_i$，即 $f_i$ 上任意第 1 型\Term{arrow}{箭头} $L_i\xrightarrow uN_i$，由 4.4.5 得到的三角形
\[\begin{tikzcd}
f_{i*}R\mathcal Hom(L_i,K_{X_i})\mathbin{\overset L\otimes}N_i\arrow[d]\arrow[dr]\\ R\mathcal Hom(N_i,f_{i*}K_{X_i})\mathbin{\overset L\otimes}N_i\arrow[r]&f_{i*}K_{X_i}
\end{tikzcd}\]
交换。事实上由\Term{functoriality}{函子性}可假设 $L_i=f_i^*N_i$，且 $u$ 由 $L_i$ 的恒等\Term{arrow}{箭头}给出。此时 $f_{i*}R\mathcal Hom(L_i,K_{X_i})\to R\mathcal Hom(N_i,f_{i*}K_{X_i})$ 是平凡\Term{duality-isomorphism}{对偶同构}的逆；取 $N_i$ 的\Term{flat-resolution}{平坦消解}与 $K_{X_i}$ 的\Term{injective-resolution}{内射消解}，即可在\Term{complex}{复形}层次直接核对三角形交换。

 b) (A) 的\Term{commutativity}{交换性}。回顾 3.7.3 的定义，只须在以下两种情形分别核对：(i) $X_i=Y_i$，$f_i$ 为 $X_i$ 的恒等\Term{arrow}{箭头}；(ii) 4.4.0 的侧面为\Term{cartesian}{笛卡尔}方形。先考虑 (i)，方形 (A) 是下图的外\Term{boundary}{边界}：
\SourcePage{117}
\[\begin{tikzcd}[column sep=small]
c'_*c'^!P\mathbin{\overset L\otimes}c''_*c''^!Q\arrow[r,"\sim"]\arrow[d,"(8)"']&c_*(c'^!P\mathbin{\overset L\otimes}_Xc''^!Q)\arrow[r]\arrow[d]&c_*c^!(P\mathbin{\overset L\otimes}Q)\arrow[d,"(9)"]\\ d'_*d'^!P\mathbin{\overset L\otimes}d''_*d''^!Q\arrow[r,"\sim"]&d_*(d'^!P\mathbin{\overset L\otimes}_Xd''^!Q)\arrow[r]&d_*d^!(P\mathbin{\overset L\otimes}Q).
\end{tikzcd}\]
左方形 $A_1$ 的水平\Term{isomorphism}{同构}为 Künneth \Term{isomorphism}{同构}；右方形 $A_2$ 的水平\Term{arrow}{箭头}由 4.2.1 给出；中间竖直\Term{arrow}{箭头}由对\Term{arrow}{箭头} $g_!(c'^!P\mathbin{\overset L\otimes}_Xc''^!Q)\to d'^!P\mathbin{\overset L\otimes}_Xd''^!Q$ 应用 $d_*$ 得到。该\Term{arrow}{箭头}通过 Künneth，由\Term{adjoint}{伴随箭头} $f'_!c'^!P\to d'^!P$、$f''_!c''^!Q\to d''^!Q$ 的\Term{external-tensor-product}{外张量积}定义。这些\Term{adjoint}{伴随箭头}与 $d'_!d'^!P$、$d''_!d''^!Q$ 的恒等\Term{arrow}{箭头}定义第 3 型\Term{arrow}{箭头}（1.4）的\Term{commutative-triangle}{交换三角形}
\[\begin{tikzcd}&c'^!P\arrow[dl]\arrow[d]&&&c''^!Q\arrow[dl]\arrow[d]\\d'_!d'^!P&d'^!P\arrow[l]&&d''_!d''^!Q&d''^!Q\arrow[l]\end{tikzcd}\]
它们分别位于
\[\begin{tikzcd}&C'\arrow[dl,"c'"']\arrow[d,"f'"]&&&C''\arrow[dl,"c''"']\arrow[d,"f''"]\\X&D'\arrow[l,"d'"]&&X&D''\arrow[l,"d''"]\end{tikzcd}\]
之上。在 $X$ 上取\Term{external-tensor-product}{外张量积}，得到\Term{commutative-triangle}{交换三角形}
\[\begin{tikzcd}&c'^!P\mathbin{\overset L\otimes}_Xc''^!Q\arrow[dl]\arrow[d]\\d'_!d'^!P\mathbin{\overset L\otimes}d''_!d''^!Q&d'^!P\mathbin{\overset L\otimes}_Xd''^!Q\arrow[l]\end{tikzcd}\qquad\text{位于}\quad\begin{tikzcd}&C\arrow[dl,"c"']\arrow[d,"g"]\\X&D\arrow[l,"d"]\end{tikzcd}\]
其\Term{commutativity}{交换性}\Term{equivalence}{等价}于把 $A_1$ 的水平\Term{arrow}{箭头}反向后所得方形的\Term{commutativity}{交换性}。另一方面还有第 3 型\Term{arrow}{箭头}的\Term{commutative-triangle}{交换三角形}
\SourcePage{118}
\[\begin{tikzcd}&c'^!P\arrow[dl]\arrow[d]&&&c''^!Q\arrow[dl]\arrow[d]\\P&d'^!P\arrow[l]&&Q&d''^!Q\arrow[l]\end{tikzcd}\]
它们由 $d'^!P,c'^!P,d''^!Q,c''^!Q$ 的恒等\Term{arrow}{箭头}定义。取\Term{external-tensor-product}{外张量积}，得到\Term{commutative-triangle}{交换三角形}
\[\begin{tikzcd}&c'^!P\mathbin{\overset L\otimes}_Xc''^!Q\arrow[dl]\arrow[d]\\P\mathbin{\overset L\otimes}Q&d'^!P\mathbin{\overset L\otimes}_Xd''^!Q\arrow[l]\end{tikzcd}\]
可将其解释为 $g:C\to D$ 上第 3 型\Term{arrow}{箭头}的\Term{morphism}{态射}
\[\begin{tikzcd}
c'^!P\mathbin{\overset L\otimes}_Xc''^!Q\arrow[r]\arrow[d]&c^!(P\mathbin{\overset L\otimes}Q)\arrow[d]\\d'^!P\mathbin{\overset L\otimes}_Xd''^!Q\arrow[r]&d^!(P\mathbin{\overset L\otimes}Q),
\end{tikzcd}\]
或 $D$ 上的交换方形
\[\begin{tikzcd}
g_*(c'^!P\mathbin{\overset L\otimes}_Xc''^!Q)\arrow[r]\arrow[d]&g_*c^!(P\mathbin{\overset L\otimes}Q)\arrow[d]\\d'^!P\mathbin{\overset L\otimes}_Xd''^!Q\arrow[r]&d^!(P\mathbin{\overset L\otimes}Q).
\end{tikzcd}\]
对后一方形应用 $d_*$ 所得的正是 $A_2$，因此 (i) 中 (A) 交换。再处理 (ii)。4.4.0 中含 $f$ 的面为\Term{cartesian}{笛卡尔}方形，故\Term{arrow}{箭头} $f'_*c'^!P\to d'^!f_*P$、$f''_*c''^!Q\to d''^!f_*Q$（3.7.3）是\Term{isomorphism}{同构}。它们的逆定义 $f',f''$ 上第 1 型\Term{arrow}{箭头} $c'^!P\to d'^!f_*P$、$c''^!Q\to d''^!f_*Q$；取\Term{direct-image}{直像}，得到第 1 型\Term{arrow}{箭头}的交换方形
\SourcePage{119}
\[\begin{tikzcd}c'_*c'^!P\arrow[d]&c'^!P\arrow[l]\arrow[d]\\ d'_*d'^!f_*P&d'^!f_*P\arrow[l]\end{tikzcd}\qquad\begin{tikzcd}c''_*c''^!Q\arrow[d]&c''^!Q\arrow[l]\arrow[d]\\ d''_*d''^!f_*Q&d''^!f_*Q\arrow[l]\end{tikzcd}\]
位于 4.4.0 中含 $f$ 的面之上。取\Term{tensor-product}{张量积}，得到 $D_{ctf}(Y)$ 中的交换方形
\[\begin{tikzcd}[column sep=small]
d'_*d'^!f_*P\mathbin{\overset L\otimes}d''_*d''^!f_*Q\arrow[r]\arrow[d]&d_*(d'^!f_*P\mathbin{\overset L\otimes}_Yd''^!f_*Q)\arrow[d]\\ f_*c'_*c'^!P\mathbin{\overset L\otimes}f_*c''_*c''^!Q\arrow[r]&d_*g_*(c'^!P\mathbin{\overset L\otimes}_Xc''^!Q).
\end{tikzcd}\tag{4.4.6}\]
左竖直\Term{arrow}{箭头}是 (8) 的逆，水平\Term{arrow}{箭头}由 Künneth 给出，右竖直\Term{arrow}{箭头}由 $Y$ 上\Term{arrow}{箭头} $f'_*c'^!P\to d'^!f_*P$、$f''_*c''^!Q\to d''^!f_*Q$ 的\Term{external-tensor-product}{外张量积}给出。同一\Term{tensor-product}{张量积}定义下图的左竖直\Term{arrow}{箭头}：
\[\begin{tikzcd}
d'^!f_*P\mathbin{\overset L\otimes}_Yd''^!f_*Q\arrow[r]\arrow[d]&d^!f_*(P\mathbin{\overset L\otimes}Q)\arrow[d]\\g_*(c'^!P\mathbin{\overset L\otimes}_Xc''^!Q)\arrow[r]&g_*c^!(P\mathbin{\overset L\otimes}Q).
\end{tikzcd}\tag{4.4.7}\]
右竖直\Term{arrow}{箭头}为 3.7.3，水平\Term{arrow}{箭头}由 4.2.1 定义。将 4.4.6 与 $d_*(4.4.7)$ 合成所得图的外\Term{boundary}{边界}，就是将 (A) 的竖直\Term{arrow}{箭头}反向所得的方形。故只须证明 4.4.7 交换，亦即 4.2.1 的\Term{arrow}{箭头}
\[c'^!P\mathbin{\overset L\otimes}_Xc''^!Q\to c^!(P\mathbin{\overset L\otimes}Q),\qquad d'^!f_*P\mathbin{\overset L\otimes}_Yd''^!f_*Q\to d^!f_*(P\mathbin{\overset L\otimes}Q)\]
定义从 $c'^!P\mathbin{\overset L\otimes}_Xc''^!Q\to d'^!f_*P\mathbin{\overset L\otimes}_Yd''^!f_*Q$ 到 $c^!(P\mathbin{\overset L\otimes}Q)\to d^!f_*(P\mathbin{\overset L\otimes}Q)$ 的\Term{morphism}{态射}，两者视为 $g$ 上第 1 型\Term{arrow}{箭头}。由\Term{adjunction}{伴随}性，归结为验证\Term{adjoint}{伴随箭头}的\Term{tensor-product}{张量积}
\[c'_*c'^!P\mathbin{\overset L\otimes}c''_*c''^!Q\to P\mathbin{\overset L\otimes}Q,\quad d'_*d'^!f_*P\mathbin{\overset L\otimes}d''_*d''^!f_*Q\to f_*P\mathbin{\overset L\otimes}f_*Q\]
在 $f$ 上第 1 型\Term{arrow}{箭头}之间，定义从 $c'_*c'^!P\to d'_*d'^!f_*P$ 与 $c''_*c''^!Q\to d''_*d''^!f_*Q$ 的\Term{tensor-product}{张量积}到 $P\to f_*P$ 与 $Q\to f_*Q$ 的\Term{tensor-product}{张量积}的\Term{morphism}{态射}。为此只须核对\Term{adjoint}{伴随箭头} $c'_*c'^!P\to P$、$d'_*d'^!f_*P\to f_*P$ 定义从 $c'_*c'^!P\to d'_*d'^!f_*P$ 到 $P\to f_*P$ 的\Term{morphism}{态射}，以及 $(c'',d'')$ 的类似断言。
\SourcePage{120}
两个情形的陈述相同，故只考虑 $(c',d')$。要证三角形
\[\begin{tikzcd}&f_*c'_*c'^!P\arrow[dl]\arrow[d,"\sim"]\\ f_*P&d'_*d'^!f_*P\arrow[l]\end{tikzcd}\]
交换，其中水平与斜\Term{arrow}{箭头}由\Term{adjoint}{伴随箭头}给出，竖直\Term{arrow}{箭头}由 3.7.3 的\Term{isomorphism}{同构}给出。按 3.7.3（SGA 4 XVIII 3.1.12.3）作为\Term{base-change}{基变换}\Term{isomorphism}{同构}的\Term{transpose}{转置}的定义，应用\Term{functor}{函子} $\operatorname{Hom}(R,-)$（$R\in\operatorname{ob}D_{ctf}(Y)$），归结为验证方形
\[\begin{tikzcd}
c'_*c'^*f^*R\arrow[r,equal]&c'_*f'^*d'^*R\\ f^*R\arrow[u]\arrow[r]&f^*d'_*d'^*R\arrow[u,"\sim"']
\end{tikzcd}\]
交换，其中左竖直\Term{arrow}{箭头}（相应地下水平\Term{arrow}{箭头}）由\Term{adjoint}{伴随箭头} $1\to c'_*c'^*$（相应地 $1\to d'_*d'^*$）定义，右竖直\Term{arrow}{箭头}由\Term{base-change}{基变换}\Term{isomorphism}{同构} $f^*d'_*\to c'_*f'^*$ 给出。依\Term{base-change}{基变换}\Term{arrow}{箭头}的定义，该方形交换。故 (A) 交换，\Term{theorem}{定理}得证。
\end{proof}
\begin{SGAresult}{推论}{4.5}
在 4.4 的假设下，有交换方形
\[\begin{tikzcd}[column sep=small]
\operatorname{Hom}(c_1'^*L_1,c_2'^!L_2)\otimes\operatorname{Hom}(c_2''^*L_2,c_1''^!L_1)\arrow[r,"4.2.5"]\arrow[d,"{3.7.6\otimes3.7.6}"']&H^0(C,K_C)\arrow[d,"g_*"]\\
\operatorname{Hom}(d_1'^*f_{1*}L_1,d_2'^!f_{2*}L_2)\otimes\operatorname{Hom}(d_2''^*f_{2*}L_2,d_1''^!f_{1*}L_1)\arrow[r,"4.2.5"]&H^0(D,K_D),
\end{tikzcd}\]
其中 $g_*$ 由\Term{adjoint}{伴随箭头} $g_*K_C=g_*g^!K_D\to K_D$ 定义。换言之，对 $u\in\operatorname{Hom}(c_1'^*L_1,c_2'^!L_2)$、$v\in\operatorname{Hom}(c_2''^*L_2,c_1''^!L_1)$，有
\SourcePage{121}
\[\langle f_*u,f_*v\rangle=g_*\langle u,v\rangle.\tag{4.5.1}\]
\end{SGAresult}
对任意\Term{proper}{固有} $S$ \Term{scheme}{概形} $t:T\to S$，以
\[\int_T:H^0(T,K_T)\longrightarrow H^0(S,K_S)=H^0(S,A)\tag{4.6}\]
记由\Term{adjoint}{伴随箭头} $t_!K_T\to A$ 定义的\Term{arrow}{箭头}。
\begin{SGAresult}{推论}{4.7（Lefschetz--Verdier 公式）}
在 4.4 的假设下，设 $Y_1=Y_2=S=D'=D''$。对 $u\in\operatorname{Hom}(c_1'^*L_1,c_2'^!L_2)$、$v\in\operatorname{Hom}(c_2''^*L_2,c_1''^!L_1)$，有
\[\langle u_*,v_*\rangle=\int_C\langle u,v\rangle,\tag{4.7.1}\]
其中 $u_*\in\operatorname{Hom}(f_{1*}L_1,f_{2*}L_2)$、$v_*\in\operatorname{Hom}(f_{2*}L_2,f_{1*}L_1)$ 是 $u,v$ 在 $S$ 上的\Term{direct-image}{直像}（3.7.6），而 $\langle u_*,v_*\rangle=\operatorname{Tr}(u_*v_*)=\operatorname{Tr}(v_*u_*)$ 是\Term{perfect-complex}{完美复形}的\Term{homomorphism}{同态} $u_*,v_*$ 在（SGA 6 I 8.3）意义下的\Term{cup-product}{杯积}（4.1.5）。
\end{SGAresult}
结合 3.6.5，特别有：
\begin{SGAresult}{推论}{4.8}
设 $t:T\to S$ 是\Term{proper}{固有} $S$ \Term{scheme}{概形}，$\Delta\subset T\times_ST$ 是\Term{diagonal}{对角线}，$F:T\to T$ 是 $S$ \Term{endomorphism}{自同态}，$T^F=(F\text{ 的图})\times_{T\times_ST}\Delta$ 是 $F$ 的\Term{fixed-point}{不动点}\Term{scheme}{概形}，$L\in\operatorname{ob}D_{ctf}(T)$，$u\in\operatorname{Hom}(F^*L,L)$ 是\Term{support}{支撑}在 $F$ 的图上的 $L$ 到 $L$ 的对应，$1\in\operatorname{Hom}(L,L)$ 是\Term{support}{支撑}在 $\Delta$ 上的恒等对应（3.2 b)）。则
\[\operatorname{Tr}(u_*)=\int_{T^F}\langle u,1\rangle,\]
其中 $u_*$ 是 $t_*L\to t_*F^*L$（\Term{inverse-image}{逆像}）与 $t_*u:t_*F^*L\to t_*L$ 复合所得的 $t_*L$ 的\Term{endomorphism}{自同态}。
\end{SGAresult}
另一方面，由 3.2 a) 与 3.6 的注记得到：
\begin{SGAresult}{推论}{4.9}
在 4.7 的假设下，设 $X_i$ 在 $S$ 上\Term{smooth}{光滑}，纯\Term{relative-dimension}{相对维数}为 $r_i$，$C'$（相应地 $C''$）是 $X$ 的闭子\Term{scheme}{概形}，
\SourcePage{122}
在 $X_2$（相应地 $X_1$）上\Term{finite}{有限}且有\Term{finite-tor-dimension}{有限 Tor 维数}。设 $\operatorname{cl}(C')\in H_{C'}^{2r_1}(X,A(r_1))$、$\operatorname{cl}(C'')\in H_{C''}^{2r_2}(X,A(r_2))$ 为其\Term{cohomology-class}{上同调类}，$\operatorname{cl}(C')_*\in\operatorname{Hom}(Rf_{1*}A,Rf_{2*}A)$、$\operatorname{cl}(C'')_*\in\operatorname{Hom}(Rf_{2*}A,Rf_{1*}A)$ 为相应\Term{homomorphism}{同态}（3.2 a)、3.7.6）。则
\[\langle\operatorname{cl}(C')_*,\operatorname{cl}(C'')_*\rangle=\int_{C'\cap C''}\operatorname{cl}(C)\operatorname{cl}(D),\]
其中
\[\operatorname{cl}(C)\operatorname{cl}(D)\in H_{C'\cap C''}^{2(r_1+r_2)}(X,A(r_1+r_2))=H^0(C'\cap C'',K_{C'\cap C''})\]
是通常的\Term{cup-product}{杯积}。
\end{SGAresult}
结合 4.3，特别有：
\begin{SGAresult}{推论}{4.10}
在 4.9 的假设下，设 $C'\cap C''$ \Term{finite}{有限}。对 $z\in C'\cap C''$，以 $(C'.C'')_z$ 记 $C'$ 与 $C''$ 在 $z$ 处的\Term{intersection-multiplicity}{交数}在 $A$ 中的像。则
\[\langle\operatorname{cl}(C')_*,\operatorname{cl}(C'')_*\rangle=\sum_{z\in C'\cap C''}(C'.C'')_z.\]
\end{SGAresult}
\textbf{4.11。} 在 4.9 的假设下，一般有
\[\langle\operatorname{cl}(C')_*,\operatorname{cl}(C'')_*\rangle=\sum_{Z\in\pi_0(C'\cap C'')}\int_Z(\operatorname{cl}(C')\operatorname{cl}(C''))_Z,\tag{4.11.1}\]
其中 $\pi_0(C'\cap C'')$ 是 $C'\cap C''$ 的\Term{connected-component}{连通分支}集，$(\operatorname{cl}(C')\operatorname{cl}(C''))_Z$ 是 $\operatorname{cl}(C')\operatorname{cl}(C'')$ 在 $H^0(Z,K_Z)$ 中的\Term{restriction}{限制}。当 $Z$ 为\Term{dimension}{维数} $\geq1$ 的分支时，由于处于“超额”情形，计算此项会遇到问题。若 $Z$ \Term{smooth}{光滑}且纯\Term{dimension}{维数}为 $d$，$C',C''$ 在 $Z$ 附近\Term{smooth}{光滑}，$X$ 是\Term{projective}{射影}的，则可证明
\[\int_Z(\operatorname{cl}(C')\operatorname{cl}(C''))_Z=\int_Zc_d(N_Z),\tag{4.11.2}\]
其中 $N_Z$ 为由\Term{normal-bundle}{法丛}\Term{exact-sequence}{正合列}
\SourcePage{123}
\[0\longrightarrow N_{Z/C'}\oplus N_{Z/C''}\longrightarrow N_{Z/X}\longrightarrow N_Z\longrightarrow0\]
定义的秩 $d$ \Term{locally-free}{局部自由}\Term{sheaf}{层}，$c_d(N_Z)\in H^{2d}(Z,A(d))$ 是其第 $d$ 个 Chern 类（SGA 5 VII）（利用 SGA 6 XIV 6.3 归结到 $K$ 理论中\Term{intersection-multiplicity}{交数}的计算，再应用 SGA 6 VII 2.5，参见 [10] 4.1；\Term{projective}{射影}性假设大概不必要）。
设 $T$ 为\Term{projective}{射影}\Term{smooth}{光滑} $S$ \Term{scheme}{概形}，取对应 $C'$ 为 $S$ \Term{endomorphism}{自同态} $F$ 的图，$C''$ 为\Term{diagonal}{对角线}，并设\Term{fixed-point}{不动点}\Term{scheme}{概形} $T^F=C'\cap C''$ \Term{smooth}{光滑}。上述\Term{sheaf}{层} $N(T^F)$ 正是 $T^F$ 的切\Term{sheaf}{层}；因此，由 Gauss--Bonnet 公式（SGA 5 VII 4.9），4.11.2 给出\Term{topology}{拓扑}学中一个著名公式的类似式：
\[\operatorname{Tr}(F,R\Gamma(T,A))=\operatorname{EP}(T^F),\tag{4.11.3}\]
其中 $\operatorname{EP}(T^F)$ 表示 $T^F_{\bar s}$ 的 \Term{euler-poincare-characteristic}{Euler--Poincaré 特征}，$\bar s$ 是 $s$ 上的\Term{geometric-point}{几何点}（同上）。据撰稿人所知，对非\Term{smooth}{光滑}且\Term{dimension}{维数} $\geq1$ 的分支 $Z$，一般尚未研究 $(\operatorname{cl}(C')\operatorname{cl}(C''))_Z$ 的计算。
\textbf{4.12。} 在 4.7 的情形下，设 $z$ 为 $C$ 的\Term{isolated-point}{孤立点}，在 $X_i$ 上\Term{projection}{投影}为 $z_i$，在 $C'$（相应地 $C''$）上\Term{projection}{投影}为 $a'$（相应地 $a''$）；设 $\bar z$ 为 $z$ 上的\Term{geometric-point}{几何点}，其在 $z_i$ 上的像为 $\bar z_i$。设 $c_2'$（相应地 $c_1''$）在 $a'$（相应地 $a''$）处\Term{quasi-finite}{拟有限}且\Term{flat}{平坦}，$L_i$ 在 $z_i$ 附近为\Term{perfect-complex}{完美复形}，$u$（相应地 $v$）在 $a'$（相应地 $a''$）附近由 $\operatorname{Tr}_{c_2'}:c'_{2!}c_2'^*L_2\to L_2$ 及 $s\in\operatorname{Hom}(c_1'^*L_1,c_2'^*L_2)$（相应地 $\operatorname{Tr}_{c_1''}:c''_{1!}c_1''^*L_1\to L_1$ 及 $t\in\operatorname{Hom}(c_2''^*L_2,c_1''^*L_1)$）给出。以 $s_{\bar z}\in\operatorname{Hom}(L_{1\bar z_1},L_{2\bar z_2})$（相应地 $t_{\bar z}\in\operatorname{Hom}(L_{2\bar z_2},L_{1\bar z_1})$）记 $s$（相应地 $t$）在\Term{fiber}{纤维}上诱导的\Term{homomorphism}{同态}。如同 3.2 a)，以 $\operatorname{cl}(c')$（相应地 $\operatorname{cl}(c'')$）记在 $a'$（相应地 $a''$）附近由 $c_2'$（相应地 $c_1''$）定义的 $X_1$ 到 $X_2$（相应地 $X_2$ 到 $X_1$）的对应，并以 $(C'.C'')_z$ 记相应\Term{local-term}{局部项} $\langle\operatorname{cl}(c'),\operatorname{cl}(c'')\rangle_z$（当 $C',C''$ 为 $X$ 的闭子\Term{scheme}{概形}，且 $X_1,X_2$ \Term{smooth}{光滑}时，这等于
\SourcePage{124}
$C',C''$ 在 $z$ 处的\Term{intersection-multiplicity}{交数}在 $A$ 中的像）。则
\[\langle u,v\rangle_{\bar z}=(C'.C'')_z\langle s_{\bar z},t_{\bar z}\rangle.\tag{4.12.1}\]
这由 4.2.6 容易得到。
\textbf{4.13。} 若 $T$ 为 $S$ 上的\Term{scheme}{概形}，以 $DF(T)$ 记具有\Term{finite}{有限}\Term{filtered-category}{滤过的} $A_T$ 模\Term{complex}{复形}的\Term{filtered}{滤过}\Term{derived-category}{导出范畴}（[9] V 1），以 $DF_{ctf}(T)$ 记由满足 $\operatorname{gr}(L)\in\operatorname{ob}D_{ctf}(T)$ 的 $L$ 组成的\Term{full-subcategory}{满子范畴}。\Term{filtered}{滤过}\Term{derive}{导出}\Term{functor}{函子}与取\Term{associated-graded}{伴随分次}相容的一系列性质，蕴含 $DF_{ctf}$ 在六运算 $f_*,f^*,f_!,f^!,\mathbin{\overset L\otimes},R\mathcal Hom$ 下\Term{stable}{稳定}。前述各条的定义和结果，作必要变更后在 $DF_{ctf}$ 中仍成立。\Term{pairing}{配对} 4.1.2 与取\Term{associated-graded}{伴随分次}相容（赋予 $K_X$ 平凡\Term{filtered}{滤过}），所以用 4.2.5 的记号，对 $L_i\in\operatorname{ob}DF_{ctf}(X_i)$、$u\in\operatorname{Hom}_{DF}(c_1^*L_1,c_2^!L_2)$、$v\in\operatorname{Hom}_{DF}(d_2^*L_2,d_1^!L_1)$，在 $H^0(E,K_E)$ 中有
\[\langle u,v\rangle=\sum_i\langle\operatorname{gr}^iu,\operatorname{gr}^iv\rangle.\tag{4.13.1}\]
正如 Deligne 所指出，这一注记特别适用于 Langlands 在 [13] 第 7 节考虑的情形。他在同上 7.11 中考虑的对应，在\Term{fixed-point}{不动点}附近保持\Term{sheaf}{层}上的\Term{finite}{有限}\Term{filtered}{滤过}，且\Term{associated-graded}{伴随分次}在\Term{fixed-point}{不动点}的像处的\Term{strict-localization}{严格局部化}上，诱导常值\Term{sheaf}{层}，或在\Term{closed-point}{闭点}的补集上的常值\Term{sheaf}{层}的\Term{extension-by-zero}{零延拓}。\Term{local-term}{局部项}的计算由 4.13 容易归结到 4.12.1，以及集中于\Term{fixed-point}{不动点}像\Term{closed-point}{闭点}上的\Term{sheaf}{层}的情形；后一情形可直接处理。应用 4.7 即得 [13] 7.11 所需的 Lefschetz 公式。指出，同上 7.12 的一般陈述也成立，其较为细致的证明将在 III B 给出。
\SourcePage{125}
\SGAsection{5}{补充}
\textbf{5.1。对应的\Term{duality}{对偶}。} 对 $i=1,2$，设 $X_i$ 为 $S$ \Term{scheme}{概形}，$X=X_1\times_SX_2$，$p_i:X\to X_i$ 为\Term{projection}{投影}，$L_i\in\operatorname{ob}D_{ctf}(X_i)$。对任意从 $L_1$ 到 $L_2$ 的\Term{cohomological-correspondence}{上同调对应} $u$，应用 $D_X$ 并作等同 $Dp_2^!\simeq p_2^*D$、$Dp_1^*\simeq p_1^!D$，得到从 $DL_2$ 到 $DL_1$ 的对应 $Du$，称为 $u$ 的\Term{duality}{对偶}。换言之，$D$ 定义\Term{homomorphism}{同态}
\[D:\operatorname{Hom}(p_1^*L_1,p_2^!L_2)\longrightarrow\operatorname{Hom}(p_2^*DL_2,p_1^!DL_1).\tag{5.1.1}\]
可验证 5.1.1 是对下列\Term{isomorphism}{同构}应用 $H^0(X,-)$ 而得：
\[D:R\mathcal Hom(p_1^*L_1,p_2^!L_2)\longrightarrow R\mathcal Hom(p_2^*DL_2,p_1^!DL_1),\tag{5.1.2}\]
该\Term{isomorphism}{同构}使方形
\[\begin{tikzcd}
p_1^*DL_1\mathbin{\overset L\otimes}p_2^*L_2\arrow[r,"3.1.1"]\arrow[d]&R\mathcal Hom(p_1^*L_1,p_2^!L_2)\arrow[d,"D"]\\p_2^*DDL_2\mathbin{\overset L\otimes}p_1^*DL_1\arrow[r,"3.1.1"]&R\mathcal Hom(p_2^*DL_2,p_1^!DL_1)
\end{tikzcd}\tag{5.1.3}\]
交换，其中左竖直\Term{arrow}{箭头}是\Term{symmetry-isomorphism}{对称同构}与\Term{biduality-isomorphism}{双对偶同构} $L_2\to DDL_2$ 定义的\Term{isomorphism}{同构}的复合。由定义容易得到以下方形交换：
\[\begin{tikzcd}
R\mathcal Hom_S(L_1,L_2)\mathbin{\overset L\otimes}R\mathcal Hom_S(L_2,L_1)\arrow[r,"4.1.3"]\arrow[d,"D\otimes D"']&K_X\arrow[d,equal]\\R\mathcal Hom_S(DL_2,DL_1)\mathbin{\overset L\otimes}R\mathcal Hom_S(DL_1,DL_2)\arrow[r,"4.1.3"]&K_X.
\end{tikzcd}\tag{5.1.4}\]
设 $c:C\to X$ 为 $S$ \Term{morphism}{态射}，置 $c_i=p_ic$；对 5.1.2 应用 $c_*c^!$（3.2.1），得到\Term{isomorphism}{同构}
\[D:\operatorname{Hom}(c_1^*L_1,c_2^!L_1)\xrightarrow\sim\operatorname{Hom}(c_2^*DL_2,c_1^!DL_1),\tag{5.1.5}\]
可验证它等同于应用\Term{functor}{函子} $D_X$ 并作等同
\SourcePage{126}
$Dc_1^*\simeq c_1^!D$、$Dc_2^!\simeq c_2^*D$ 所得的\Term{isomorphism}{同构}。若 $c':C'\to X$、$c'':C''\to X$ 是 $S$ \Term{morphism}{态射}，则由 5.1.4 得，对 $u\in\operatorname{Hom}(c_1'^*L_1,c_2'^!L_2)$、$v\in\operatorname{Hom}(c_2''^*L_2,c_1''^!L_1)$，有
\[\langle u,v\rangle=\langle Du,Dv\rangle,\tag{5.1.6}\]
记号如 4.2.5。另一方面，可验证在 3.3 的情形下，方形
\[\begin{tikzcd}
f_*R\mathcal Hom_S(L_1,L_2)\arrow[r,"3.3.1"]\arrow[d,"f_*D"']&R\mathcal Hom_S(f_{1!}L_1,f_{2*}L_2)\arrow[d,"D"]\\ f_*R\mathcal Hom_S(DL_2,DL_1)\arrow[r,"3.3.1"]&R\mathcal Hom_S(f_{2!}DL_2,f_{1*}DL_1)
\end{tikzcd}\tag{5.1.7}\]
交换，因此在 3.7.6 的情形下有
\[f_*Du=Df_*u.\tag{5.1.8}\]
\textbf{5.2。\Term{composition-of-correspondences}{对应的复合}。} 对 $i=1,2,3$，设 $X_i$ 为 $S$ \Term{scheme}{概形}，$X_{ij}=X_i\times_SX_j$，$X=X_1\times_SX_2\times_SX_3$，$p_i:X\to X_i$、$p_{ij}:X\to X_{ij}$ 为\Term{projection}{投影}，$L_i\in\operatorname{ob}D_{ctf}(X_i)$。有 $X=X_{12}\times_{X_2}X_{23}$，以及
\[(DL_1\mathbin{\overset L\otimes}_SL_2)\mathbin{\overset L\otimes}_{X_2}(DL_2\mathbin{\overset L\otimes}_SL_3)=p_1^*DL_1\mathbin{\overset L\otimes}p_2^*L_2\mathbin{\overset L\otimes}p_2^*DL_2\mathbin{\overset L\otimes}p_3^*L_3.\]
故\Term{evaluation-map}{求值箭头} $L_2\mathbin{\overset L\otimes}DL_2\to K_{X_2}$ 定义\Term{pairing}{配对}
\[(DL_1\mathbin{\overset L\otimes}_SL_2)\mathbin{\overset L\otimes}_{X_2}(DL_2\mathbin{\overset L\otimes}_SL_3)\longrightarrow p_{13}^!(DL_1\mathbin{\overset L\otimes}_SL_3),\tag{5.2.1}\]
其中右端由 Künneth（1.7.3）等同于 $(DL_1\mathbin{\overset L\otimes}_SL_3)\mathbin{\overset L\otimes}_SK_{X_2}$。结合 3.1.1、3.1.2，可写成
\[R\mathcal Hom_S(L_1,L_2)\mathbin{\overset L\otimes}_{X_2}R\mathcal Hom_S(L_2,L_3)\longrightarrow p_{13}^!R\mathcal Hom_S(L_1,L_3)=R\mathcal Hom(p_1^*L_1,p_3^!L_3),\tag{5.2.2}\]
\SourcePage{127}
右端的等同由诱导\Term{isomorphism}{同构}（SGA 4 XVIII 3.1.12.2）给出。因此得到称为\Term{composition-of-correspondences}{对应的复合}的\Term{pairing}{配对}
\[\operatorname{Hom}_S(L_1,L_2)\otimes\operatorname{Hom}_S(L_2,L_3)\longrightarrow\operatorname{Hom}(p_1^*L_1,p_3^!L_3),\tag{5.2.3}\]
记为 $(u,v)\mapsto vu$。若 $p_{ij,i}:X_{ij}\to X_i$、$p_{ij,j}:X_{ij}\to X_j$ 表示\Term{projection}{投影}，可验证 5.2.2 也可由以下\Term{arrow}{箭头}的复合得到：
\[\begin{split}
R\mathcal Hom_S(L_1,L_2)\mathbin{\overset L\otimes}_{X_2}R\mathcal Hom_S(L_2,L_3)&\longrightarrow R\mathcal Hom(p_1^*L_1,p_{12}^*p_{12,2}^!L_2)\\&\mathbin{\overset L\otimes}R\mathcal Hom(p_{23}^!p_{23,2}^*L_2,p_3^!L_3),
\end{split}\]
（这是\Term{canonical}{典范}\Term{arrow}{箭头} $p_{12}^*R\mathcal Hom_S(L_1,L_2)\to R\mathcal Hom(p_1^*L_1,p_{12}^*p_{12,2}^!L_2)$、$p_{23}^*R\mathcal Hom_S(L_2,L_3)\to R\mathcal Hom(p_{23}^!p_{23,2}^*L_2,p_3^!L_3)$ 的\Term{tensor-product}{张量积}），再用\Term{base-change}{基变换}\Term{arrow}{箭头} $p_{12}^*p_{12,2}^!L_2\to p_{23}^!p_{23,2}^*L_2$ 所定义的\Term{arrow}{箭头}，以及 $R\mathcal Hom$ 上通常的复合\Term{arrow}{箭头}。
设 $c':C'\to X_{12}$、$c'':C''\to X_{23}$ 为对应，$c:C=C'\times_{X_2}C''\to X$ 为“复合”对应。由 5.2.2（或以与上述类似的复合直接定义），得到\Term{pairing}{配对}
\[R\mathcal Hom(c_1'^*L_1,c_2'^!L_2)\mathbin{\overset L\otimes}_{X_2}R\mathcal Hom(c_2''^*L_2,c_3''^!L_3)\longrightarrow R\mathcal Hom(c_1^*L_1,c_3^!L_3),\tag{5.2.4}\]
从而有“复合”\Term{pairing}{配对} $(u,v)\mapsto vu$：
\[\operatorname{Hom}(c_1'^*L_1,c_2'^!L_2)\otimes\operatorname{Hom}(c_2''^*L_2,c_3''^!L_3)\longrightarrow\operatorname{Hom}(c_1^*L_1,c_3^!L_3).\tag{5.2.5}\]
可验证\Term{composition-of-correspondences}{对应的复合}与\Term{direct-image}{直像}相容，意义如下。对 $1\leq i\leq3$，设 $f_i:X_i\to Y_i$ 为 $S$ \Term{morphism}{态射}，$f_2$ \Term{proper}{固有}；$c',c'',c$ 如上，$c',c''$ \Term{proper}{固有}；$d':D'\to Y_{12}$、$d'':D''\to Y_{23}$ 为\Term{proper}{固有}\Term{morphism}{态射}，$d:D=D'\times_{Y_2}D''\to Y$，并有交换方形
\SourcePage{128}
\[\begin{tikzcd}X_{12}\arrow[d,"f_{12}"']&C'\arrow[l]\arrow[d]\\Y_{12}&D'\arrow[l]\end{tikzcd}\qquad\begin{tikzcd}X_{23}\arrow[d,"f_{23}"']&C''\arrow[l]\arrow[d]\\Y_{23}&D''\arrow[l]\end{tikzcd}.\]
于是有交换方形
\[\begin{tikzcd}[column sep=small]
\operatorname{Hom}(c_1'^*L_1,c_2'^!L_2)\otimes\operatorname{Hom}(c_2''^*L_2,c_3''^!L_3)\arrow[r,"5.2.5"]\arrow[d,"{f_{12*}\otimes f_{23*}}"']&\operatorname{Hom}(c_1^*L_1,c_3^!L_3)\arrow[d,"f_{13*}"]\\
\operatorname{Hom}(d_1'^*M_1,d_2'^!M_2)\otimes\operatorname{Hom}(d_2''^*M_2,d_3''^!M_3)\arrow[r,"5.2.5"]&\operatorname{Hom}(d_1^*M_1,d_3^!M_3),
\end{tikzcd}\tag{5.2.6}\]
其中 $M_1=f_{1!}L_1$、$M_2=f_{2*}L_2$、$M_3=f_{3*}L_3$；也有以 $R\mathcal Hom$ 代替的类似交换方形，留给读者写出。
最后设 $X_1=X_3$、$L_1=L_3$。基本\Term{pairing}{配对} 4.1.2 \Term{solvable}{可解}释为以下两个\Term{pairing}{配对}的复合：
\[q_1^*DL_1\mathbin{\overset L\otimes}q_2^*L_2\mathbin{\overset L\otimes}q_2^*DL_2\mathbin{\overset L\otimes}q_1^*L_1\longrightarrow q_1^*DL_1\mathbin{\overset L\otimes}q_2^*K_{X_2}\mathbin{\overset L\otimes}q_1^*L_1,\tag{5.2.7}\]
它由 $L_2\mathbin{\overset L\otimes}DL_2\to K_{X_2}$ 定义；\Term{equivalence}{等价}地，它是经\Term{diagonal}{对角线} $X_{12}\to X_1\times_SX_2\times_SX_1$ 由 5.2.1 在 $X_{12}$ 上诱导的\Term{pairing}{配对}（$q_i:X_{12}\to X_i$ 为\Term{projection}{投影}），以及由 $DL_1\mathbin{\overset L\otimes}L_1\to K_{X_1}$ 定义的
\[q_1^*DL_1\mathbin{\overset L\otimes}q_2^*K_{X_2}\mathbin{\overset L\otimes}q_1^*L_1\longrightarrow K_{X_{12}}.\tag{5.2.8}\]
因此对 $u\in\operatorname{Hom}_S(L_1,L_2)$、$v\in\operatorname{Hom}_S(L_2,L_1)$，在 $H^0(X_1\times_SX_2,K_{X_1\times_SX_2})$ 中有
\[\langle u,v\rangle=\langle vu,1\rangle,\tag{5.2.9}\]
其中 $1$ 是 $X_1\times_SX_1$ 上\Term{support}{支撑}在\Term{diagonal}{对角线}的 $L_1$ 到 $L_1$ 的恒等对应（3.2 b)）；右端\Term{cup-product}{杯积}取 4.2.5 的意义，注意 $X\times_{X_1\times_SX_1}X_1$（由 $p_{13}:X\to X_{13}=X_1\times_SX_1$ 与\Term{diagonal}{对角线} $X_1\to X_1\times_SX_1$ 取\Term{fiber-product}{纤维积}）等同于 $X_1\times_SX_2$。同样，更一般地，若 $c':C'\to X_{12}$、$c'':C''\to X_{21}$ 为对应，$c:C=C'\times_{X_2}C''\to X_1\times_SX_2\times_SX_1$ 为复合对应，
\SourcePage{129}
$d:C'\times_{X_{12}}C''=D\to X_{12}$ 为交对应，则对 $u\in\operatorname{Hom}(c_1'^*L_1,c_2'^!L_2)$、$v\in\operatorname{Hom}(c_2''^*L_2,c_1''^!L_1)$，在 $H^0(D,K_D)$ 中有
\[\langle u,v\rangle=\langle vu,1\rangle,\tag{5.2.10}\]
其中 $1$ 如上为\Term{diagonal-correspondence}{对角对应}，$D$ 等同于 $C\to X_1\times_SX_1$ 与\Term{diagonal}{对角线} $X_1\to X_1\times_SX_1$ 的\Term{fiber-product}{纤维积}。5.2.9、5.2.10 推广了\Term{constant-coefficients}{常系数}情形的一个熟知公式（[12] 1.3.6 (ii) c)）。
\textbf{5.3。对应的\Term{external-tensor-product}{外张量积}。} 对 $1\leq i\leq4$，设 $X_i$ 为 $S$ \Term{scheme}{概形}，$X=X_1\times_SX_2\times_SX_3\times_SX_4$，$X_{ij}=X_i\times_SX_j$，$L_i\in\operatorname{ob}D_{ctf}(X_i)$。Künneth \Term{isomorphism}{同构}（1.7.6、2.2.4）$DL_1\mathbin{\overset L\otimes}_SDL_3\xrightarrow\sim D(L_1\mathbin{\overset L\otimes}_SL_3)$ 在 $X$ 上定义\Term{isomorphism}{同构}
\[(DL_1\mathbin{\overset L\otimes}_SL_2)\mathbin{\overset L\otimes}_S(DL_3\mathbin{\overset L\otimes}_SL_4)\xrightarrow\sim D(L_1\mathbin{\overset L\otimes}_SL_3)\mathbin{\overset L\otimes}_S(L_2\mathbin{\overset L\otimes}_SL_4),\tag{5.3.1}\]
由 3.1.1 可写为
\[R\mathcal Hom_S(L_1,L_2)\mathbin{\overset L\otimes}_SR\mathcal Hom_S(L_3,L_4)\xrightarrow\sim R\mathcal Hom_S(L_1\mathbin{\overset L\otimes}_SL_3,L_2\mathbin{\overset L\otimes}_SL_4).\tag{5.3.2}\]
于是得到称为\Term{external-tensor-product}{外张量积}的\Term{pairing}{配对}
\[\operatorname{Hom}_S(L_1,L_2)\otimes\operatorname{Hom}_S(L_3,L_4)\longrightarrow\operatorname{Hom}_S(L_1\mathbin{\overset L\otimes}_SL_3,L_2\mathbin{\overset L\otimes}_SL_4),\tag{5.3.3}\]
记为 $(u,v)\mapsto u\mathbin{\overset L\otimes}_Sv$。\Term{pairing}{配对} 5.3.2 也\Term{solvable}{可解}释为\Term{external-tensor-product}{外张量积}
\[R\mathcal Hom(E_1,E_2)\mathbin{\overset L\otimes}_SR\mathcal Hom(E_3,E_4)\xrightarrow\sim R\mathcal Hom(E_1\mathbin{\overset L\otimes}_SE_3,E_2\mathbin{\overset L\otimes}_SE_4)\]
（此处 $E_1=p_{12,1}^*L_1$ 等，记号的意义显然）与 Künneth \Term{arrow}{箭头} $E_2\mathbin{\overset L\otimes}_SE_4\to p_{24}^!(L_2\mathbin{\overset L\otimes}_SL_4)$（1.7.3）所定义的\Term{isomorphism}{同构}的复合。
\SourcePage{130}
设 $c':C'\to X_{12}$、$c'':C''\to X_{34}$ 为对应，$c:C=C'\times_SC''\to X=X_{13}\times_SX_{24}$ 为乘积对应。由 5.3.2 得到（或以类似方式定义）\Term{pairing}{配对}
\[\begin{split}
R\mathcal Hom(c_1'^*L_1,c_2'^!L_2)\mathbin{\overset L\otimes}_SR\mathcal Hom(c_3''^*L_3,c_4''^!L_4)\longrightarrow R\mathcal Hom(c_{13}^*(L_1\mathbin{\overset L\otimes}_SL_3),c_{24}^!(L_2\mathbin{\overset L\otimes}_SL_4)),
\end{split}\tag{5.3.4}\]
从而有\Term{external-tensor-product}{外张量积}\Term{pairing}{配对} $(u,v)\mapsto u\mathbin{\overset L\otimes}_Sv$：
\[\begin{split}
\operatorname{Hom}(c_1'^*L_1,c_2'^!L_2)\otimes\operatorname{Hom}(c_3''^*L_3,c_4''^!L_4)\longrightarrow\operatorname{Hom}(c_{13}^*(L_1\mathbin{\overset L\otimes}_SL_3),c_{24}^!(L_2\mathbin{\overset L\otimes}_SL_4)).
\end{split}\tag{5.3.5}\]
可验证这些\Term{pairing}{配对}与\Term{direct-image}{直像}相容，意义如下。对 $1\leq i\leq4$，设 $f_i:X_i\to Y_i$ 为 $S$ \Term{morphism}{态射}；$c',c'',c$ 如上，$c',c''$ \Term{proper}{固有}；$d':D'\to Y_{12}$、$d'':D''\to Y_{34}$ 为\Term{proper}{固有}\Term{morphism}{态射}，$d:D=D'\times_SD''\to Y=Y_{12}\times_SY_{34}$，并有交换方形
\[\begin{tikzcd}X_{12}\arrow[d,"f_{12}"']&C'\arrow[l]\arrow[d]\\Y_{12}&D'\arrow[l]\end{tikzcd}\qquad\begin{tikzcd}X_{34}\arrow[d,"f_{34}"']&C''\arrow[l]\arrow[d]\\Y_{34}&D''\arrow[l]\end{tikzcd}.\]
则有交换方形
\[\begin{tikzcd}[column sep=small]
\operatorname{Hom}(c_1'^*L_1,c_2'^!L_2)\otimes\operatorname{Hom}(c_3''^*L_3,c_4''^!L_4)\arrow[r,"5.3.5"]\arrow[d,"{f_{12*}\otimes f_{34*}}"']&\operatorname{Hom}(c_{13}^*L_{13},c_{24}^!L_{24})\arrow[d,"{f_{13,24*}}"]\\
\operatorname{Hom}(d_1'^*M_1,d_2'^!M_2)\otimes\operatorname{Hom}(d_3''^*M_3,d_4''^!M_4)\arrow[r,"5.3.5"]&\operatorname{Hom}(d_{13}^*M_{13},d_{24}^!M_{24}),
\end{tikzcd}\tag{5.3.6}\]
其中 $i=1,3$ 时 $M_i=f_{i!}L_i$，否则 $M_i=f_{i*}L_i$；$L_{ij}=L_i\mathbin{\overset L\otimes}_SL_j$、$M_{ij}=M_i\mathbin{\overset L\otimes}_SM_j$（也有用 $R\mathcal Hom$ 写出的类似方形）。
\SourcePage{131}
设 $X_3=X_2$、$X_4=X_1$、$L_3=L_2$、$L_4=L_1$。与 5.2 最后一段类似的推理表明，对 $u\in\operatorname{Hom}_S(L_1,L_2)$、$v\in\operatorname{Hom}_S(L_2,L_1)$，有
\[\langle u,v\rangle=\langle u\mathbin{\overset L\otimes}_Sv,1\rangle,\tag{5.3.7}\]
其中 $1$ 是 $X_{12}$ 上 $L_{12}$ 到 $L_{12}$ 的\Term{diagonal-correspondence}{对角对应}，右端\Term{cup-product}{杯积}取 4.2.5 的意义。更一般地，若 $c':C'\to X_{12}$、$c'':C''\to X_{21}$ 为对应，它们的“交”$C=C'\times_{X_{12}}C''$ 等同于 $C'\times_SC''\to X=X_{12}\times_SX_{21}$ 与\Term{diagonal}{对角线} $X_{12}\to X$ 的\Term{fiber-product}{纤维积}；对 $u\in\operatorname{Hom}(c_1'^*L_1,c_2'^!L_2)$、$v\in\operatorname{Hom}(c_2''^*L_2,c_1''^!L_1)$，在 $H^0(C,K_C)$ 中有
\[\langle u,v\rangle=\langle u\mathbin{\overset L\otimes}_Sv,1\rangle,\tag{5.3.8}\]
其中 $1$ 如上表示\Term{diagonal-correspondence}{对角对应}。因此有两种“对角”方法（5.2.10、5.3.8）来计算“Verdier \Term{local-term}{局部项}”$\langle u,v\rangle$。
\SourcePage{132}
\SGAsection{6}{附录：凝聚层的 Lefschetz 公式}
本节中，$S$ 表示 Noether \Term{scheme}{概形}。所考虑的 $S$ \Term{scheme}{概形}均假设\Term{separated}{分离}且\Term{finite-type}{有限型}。回顾 Nagata 的结果：这些\Term{scheme}{概形}之间的任意 $S$ \Term{morphism}{态射}均\Term{compactifiable}{可紧化}，即\Term{decomposable}{可分解}为\Term{open-immersion}{开浸入}随后接一个\Term{proper}{固有}\Term{morphism}{态射}。对任意 $S$ \Term{scheme}{概形} $X$，以 $D(X)$ 记 $\mathcal O_X$ \Term{module-category}{模范畴}的\Term{derived-category}{导出范畴}。
\textbf{6.1。} 回顾（[8] 附录），对任意 $S$ \Term{morphism}{态射} $f:X\to Y$，有\Term{functor}{函子} $f^!:D^+(Y)\to D^+(X)$（对复合 $gf$ 有传递\Term{isomorphism}{同构} $(gf)^!=f^!g^!$，满足\Term{cocycle-condition}{余循环条件}）。当 $f$ \Term{proper}{固有}时，$f^!$ 是\Term{functor}{函子} $f_*:D^+(X)_{coh}\to D^+(Y)_{coh}$ 的右“\Term{adjunction}{伴随}”（下标 $coh$ 表示由\Term{cohomology}{上同调}\Term{coherent}{凝聚}的\Term{complex}{复形}组成的\Term{full-subcategory}{满子范畴}）；此\Term{adjunction}{伴随}具体为\Term{isomorphism}{同构}
\[f_*R\mathcal Hom(E,f^!F)\xrightarrow\sim R\mathcal Hom(f_*E,F),\tag{6.1.1}\]
其中 $E\in\operatorname{ob}D(X)_{coh}$、$F\in\operatorname{ob}D^+(Y)$；定义方式类似（SGA 4 XVIII 3.1.9.6）。[8] 附录 5 的讨论连同\Term{projective-space}{射影空间}的\Term{duality}{对偶}\Term{theorem}{定理}表明：若 $f$ \Term{smooth}{光滑}且纯\Term{dimension}{维数}为 $d$，有\Term{canonical}{典范}\Term{isomorphism}{同构}
\[f^!L\xrightarrow\sim f^*L\otimes\Omega^d_{X/Y}[d];\tag{6.1.2}\]
若 $f$ 还\Term{proper}{固有}，则由\Term{adjoint}{伴随箭头} $f_*f^!\to1$ 经 6.1.2 得到的\Term{arrow}{箭头}
\[\operatorname{Tr}_f:f_*(f^*L\otimes\Omega^d_{X/Y}[d])\longrightarrow L\tag{6.1.3}\]
与\Term{trace-morphism}{迹态射}（[8] VI）一致。
\textbf{6.2。} 与\Term{discrete}{离散}系数情形不同，此处没有在“六运算”$(f_*,f^*,f^!,f_!,\mathbin{\overset L\otimes},R\mathcal Hom)$ 下\Term{stable}{稳定}的合理的 $D$ 的\Term{subcategory}{子范畴}。最自然的\Term{finiteness}{有限性}概念似乎是相对\Term{perfect}{完美}性（SGA 6 III 4）。若 $X$ 是 $S$ \Term{scheme}{概形}，$L\in\operatorname{ob}D(X)$，
\SourcePage{133}
回顾称 $L$ 相对于 $S$ \Term{perfect}{完美}，是指 $L$ 相对于 $S$ \Term{finite-tor-dimension}{有限 Tor 维数}且\Term{cohomology}{上同调}\Term{coherent}{凝聚}。以 $D(X)_{parf/S}$ 记由相对于 $S$ \Term{perfect}{完美}的\Term{complex}{复形}组成的 $D(X)$ 的\Term{full-subcategory}{满子范畴}。它一般在 $\mathbin{\overset L\otimes}$ 和 $R\mathcal Hom$ 下均不\Term{stable}{稳定}。不过，若 $E$ 为（绝对意义下的）\Term{perfect-complex}{完美复形}，$F$ 相对于 $S$ \Term{perfect}{完美}，则 $R\mathcal Hom(E,F)$ 相对于 $S$ \Term{perfect}{完美}（SGA 6 III 4.9）。另一方面，对 $i=1,2$，设 $X_i$ 为 $S$ \Term{scheme}{概形}，$L_i\in\operatorname{ob}D(X_i)$；如同 1.5，置
\[L_1\mathbin{\overset L\otimes}_SL_2=Lp_1^*L_1\mathbin{\overset L\otimes}Lp_2^*L_2,\]
其中 $p_i:X_1\times_SX_2\to X_i$ 为\Term{projection}{投影}。只有 $X_1,X_2$ 在 $S$ 上 \Term{tor-independence}{Tor 无关}，即（SGA 6 III 1.5）$\operatorname{Tor}_i^S(\mathcal O_X,\mathcal O_Y)=0$（$i\ne0$）时，考虑这一\Term{object}{对象}才合理。在此假设下，若 $L_1,L_2$ 相对于 $S$ \Term{perfect}{完美}，则 $L_1\mathbin{\overset L\otimes}_SL_2$ 也如此（SGA 6 III 4.7）。
设 $f:X\to Y$ 为 $S$ \Term{morphism}{态射}。若 $f$ \Term{proper}{固有}，则 $f_*D(X)_{parf/S}\subset D(Y)_{parf/S}$（SGA 6 III 4.8）；若 $f$ \Term{finite-tor-dimension}{有限 Tor 维数}，则 $f^*D(Y)_{parf/S}\subset D(X)_{parf/S}$（SGA 6 III 4.5），且 $f^!D(Y)_{parf/S}\subset D(X)_{parf/S}$（SGA 6 III 4.9）。对任意 $S$ \Term{scheme}{概形} $a:X\to S$，置
\[K_X=a^!\mathcal O_S,\tag{6.2.1}\]
并以 $D_X$（或简记 $D$）记\Term{functor}{函子} $R\mathcal Hom(-,K_X)$。若 $X$ 在 $S$ 上有\Term{finite-tor-dimension}{有限 Tor 维数}，则 $K_X$ 相对于 $S$ \Term{perfect}{完美}；对任意 $L\in\operatorname{ob}D(X)_{parf/S}$，有 $DL\in\operatorname{ob}D(X)_{parf/S}$，且\Term{canonical}{典范}\Term{arrow}{箭头} $L\to DDL$ 是\Term{isomorphism}{同构}（SGA 6 III 4.9）。
设 $f:X\to Y$ 为 $S$ \Term{morphism}{态射}，$L\in\operatorname{ob}D(X)$、$M\in\operatorname{ob}D(Y)$。若 $f^*M$ 有定义（例如 $M\in\operatorname{ob}D^-(Y)$，或 $f$ \Term{finite-tor-dimension}{有限 Tor 维数}），置
\[\operatorname{Hom}_f^{(1)}(L,M)=\operatorname{Hom}(f^*M,L),\qquad\operatorname{Hom}_f^{(2)}(L,M)=\operatorname{Hom}(L,f^*M).\]
若 $f^!M$ 有定义（例如 $M\in\operatorname{ob}D^+(Y)$，或 $f$ \Term{smooth}{光滑}，或 $f$ \Term{finite}{有限}且\Term{finite-tor-dimension}{有限 Tor 维数}，或 $f$ 为此类\Term{morphism}{态射}的复合；这些情形由 6.1.1、6.1.2，$f^!$ 可延拓为
\SourcePage{134}
整个 $D(Y)$ 上的\Term{functor}{函子}），置
\[\operatorname{Hom}_f^{(3)}(L,M)=\operatorname{Hom}(L,f^!M),\qquad\operatorname{Hom}_f^{(4)}(L,M)=\operatorname{Hom}(f^!M,L).\]
如同 1.4，称 $\operatorname{Hom}_f^{(i)}(L,M)$ 的元素为 $f$ 上从 $L$ 到 $M$ 的第 $i$ 型\Term{arrow}{箭头}（或从 $L$ 到 $M$ 的第 $i$ 型 $f$ \Term{arrow}{箭头}）。对固定的 $i$，其复合及相应的\Term{fibred-category}{纤维范畴}与\Term{cofibred-category}{上纤维范畴}留给读者具体写出。这里只回顾部分\Term{adjunction-formula}{伴随公式}：
\[\operatorname{Hom}_f^{(1)}(L,M)=\operatorname{Hom}(M,f_*L)\quad(M\in\operatorname{ob}D^-(Y),\ L\in\operatorname{ob}D^+(X)),\]
\[\operatorname{Hom}_f^{(3)}(L,M)=\operatorname{Hom}(f_*L,M)\quad(f\text{ 固有},\ L\in\operatorname{ob}D(X)_{coh},\ M\in\operatorname{ob}D^+(Y)).\]
\textbf{6.3。} 在 1.6 的情形下，设 $X_1,X_2$ 在 $S$ 上 \Term{tor-independence}{Tor 无关}，$Y_1,Y_2$ 也如此。对 $L_i\in\operatorname{ob}D(X_i)_{parf/S}$，Künneth \Term{arrow}{箭头} 1.6.4 有定义（利用相对\Term{finite-tor-dimension}{有限 Tor 维数}在\Term{direct-image}{直像}下保持这一事实，SGA 6 III 3.7.1），且为\Term{isomorphism}{同构}。验证时可假设 $S,Y_i$ \Term{affine}{仿射}，再经 Čech 型计算归结到 $X_1,X_2$ 也\Term{affine}{仿射}的情形，此时结论由 \Term{tor-independence}{Tor 无关}假设平凡地得到。另一方面，对 $M_i\in\operatorname{ob}D(Y_i)_{parf/S}$，若 $f_1^!M_1$ 或 $f_2^!M_2$ 属于 $D^b$，则 Künneth \Term{arrow}{箭头} 1.7.3 有定义\footnote{在\Term{proper}{固有}或可\Term{smooth}{光滑}化的情形，\Term{arrow}{箭头}的定义是直接的；一般情形可用\Term{cohomological-descent}{上同调下降}归结到可\Term{smooth}{光滑}化情形。}；若 $f_1,f_2$ \Term{finite-tor-dimension}{有限 Tor 维数}，则它为\Term{isomorphism}{同构}。这由上述 1.6.4 型 Künneth \Term{isomorphism}{同构}结合 6.1.2，以及 $f_i$ 为\Term{closed-immersion}{闭浸入}时的 6.1.1 容易验证。特别地，若 $X_1,X_2$ 在 $S$ 上\Term{finite-tor-dimension}{有限 Tor 维数}且 \Term{tor-independence}{Tor 无关}，1.7.6 为\Term{isomorphism}{同构}。
如同 2.1，对 $(u,v)$ 为 $(2,1)$ 型（相应地 $(3,4)$ 型，并要求 $f$ \Term{proper}{固有}，以避免 $f_!$ 所导致的\Term{projective-object}{投射对象}），可定义 $R\mathcal Hom(u,v)$。在适当的\Term{degree}{次数}假设下（例如见 [9] V 2.3），\Term{evaluation-map}{求值箭头}（2.2.2）和\Term{tensor-product}{张量积}\Term{arrow}{箭头}（2.2.3、2.2.4）有定义；且只要 $E_1,E_2$ 为（绝对意义下的）\Term{perfect-complex}{完美复形}，
\SourcePage{135}
2.2.3、2.2.4（在其有定义的前提下）便是\Term{isomorphism}{同构}。
在 2.4 的情形中，设 $(X_1,X_2)$、$(X_1,Y_2)$、$(Y_1,X_2)$、$(Y_1,Y_2)$ 均在 $S$ 上 \Term{tor-independence}{Tor 无关}，$f_1$ \Term{proper}{固有}，$E_i,F_i,P_i,Q_i,R\mathcal Hom(E_i,F_i),R\mathcal Hom(P_i,Q_i)$ 相对于 $S$ \Term{perfect}{完美}，且 Künneth \Term{arrow}{箭头}
\[f_1^!Q_1\mathbin{\overset L\otimes}_SQ_2\longrightarrow(f_1\times_SY_2)^!(Q_1\mathbin{\overset L\otimes}_SQ_2),\quad f_1^!Q_1\mathbin{\overset L\otimes}_SF_2\longrightarrow(f_1\times_SY_2)^!(Q_1\mathbin{\overset L\otimes}_SF_2)\]
为\Term{isomorphism}{同构}。则 2.4.1 与 2.5 的图有定义，后者交换（但水平\Term{arrow}{箭头}未必为\Term{isomorphism}{同构}）。验证与 2.5 相同。
\textbf{6.4。} 对 $i=1,2$，设 $X_i$ 为 $S$ \Term{scheme}{概形}，$X=X_1\times_SX_2$，$p_i:X\to X_i$、$q_i:X_i\to S$ 为\Term{canonical}{典范}\Term{projection}{投影}。设 $X_1,X_2$ \Term{tor-independence}{Tor 无关}，$q_1$ \Term{finite-tor-dimension}{有限 Tor 维数}。再设 $L_i\in\operatorname{ob}D(X_i)_{parf/S}$。由 2.2.4 的类似式，有\Term{canonical}{典范}\Term{arrow}{箭头}
\[DL_1\mathbin{\overset L\otimes}_SL_2\longrightarrow R\mathcal Hom(L_1\mathbin{\overset L\otimes}_S\mathcal O_{X_2},q_1^!\mathcal O_S\mathbin{\overset L\otimes}_SL_2),\tag{6.4.1}\]
与 Künneth \Term{isomorphism}{同构}（6.3）$q_1^!\mathcal O_S\mathbin{\overset L\otimes}_SL_2\xrightarrow\sim p_2^!L_2$ 复合，得到\Term{arrow}{箭头}
\[DL_1\mathbin{\overset L\otimes}_SL_2\longrightarrow R\mathcal Hom(p_1^*L_1,p_2^!L_2).\tag{6.4.2}\]
6.4.1（因而 6.4.2）是\Term{isomorphism}{同构}，归结到 $q_1$ 为\Term{closed-immersion}{闭浸入}的情形即可看出。如同 3.1.2，置
\[R\mathcal Hom(p_1^*L_1,p_2^!L_2)=R\mathcal Hom_S(L_1,L_2),\qquad\operatorname{Hom}(p_1^*L_1,p_2^!L_2)=\operatorname{Hom}_S(L_1,L_2),\tag{6.4.3}\]
并称 $\operatorname{Hom}_S(L_1,L_2)$ 的元素为从 $L_1$ 到 $L_2$ 的\Term{cohomological-correspondence}{上同调对应}。若 $c:C\to X$ 为对应（应用中为\Term{proper}{固有}），则由\Term{induction-formula}{诱导公式}（[8] III 8.8）\footnote{所引文献的“可\Term{smooth}{光滑}化”假设不必要。}有
\SourcePage{136}
\[c^!R\mathcal Hom(p_1^*L_1,p_2^!L_2)=R\mathcal Hom(c_1^*L_1,c_2^!L_2),\tag{6.4.4}\]
称 $\operatorname{Hom}(c_1^*L_1,c_2^!L_2)=H^0(X,c_*c^!R\mathcal Hom(p_1^*L_1,p_2^!L_2))$ 的元素为\Term{support}{支撑}在 $C$ 中的 $L_1$ 到 $L_2$ 的（\Term{cohomological-correspondence}{上同调}）对应。后面将给出类似 3.2 的例子。
\textbf{6.5。} 对 $i=1,2$，设 $f_i:X_i\to Y_i$ 为\Term{proper}{固有} $S$ \Term{morphism}{态射}，$f=f_1\times_Sf_2:X\to Y$。设 $X_i,Y_i$ 在 $S$ 上有\Term{finite-tor-dimension}{有限 Tor 维数}，且 2.4.0 的方形 \Term{tor-independence}{Tor 无关}。设 $L_i\in\operatorname{ob}D(X_i)_{parf/S}$。如同 3.3.1，定义\Term{isomorphism}{同构}
\[f_*R\mathcal Hom_S(L_1,L_2)\xrightarrow\sim R\mathcal Hom_S(f_{1*}L_1,f_{2*}L_2),\tag{6.5.1}\]
应用 $R\Gamma(Y,-)$，得到\Term{isomorphism}{同构}
\[f_*:\operatorname{Hom}_S(L_1,L_2)\xrightarrow\sim\operatorname{Hom}_S(f_{1*}L_1,f_{2*}L_2);\tag{6.5.2}\]
当 $Y_1=Y_2=S$ 时，简记为 $u\mapsto u_*$。由 6.3 所述 2.5 的类似式，有与 3.4.1 类似的\Term{commutative-diagram}{交换图}。从而得到类似 6.5 的推论；若 $Y_1=Y_2=S$ 是域的谱，则 $u_*$ 在\Term{cohomology}{上同调}上的\Term{action}{作用}有类似 3.6 的描述。
给定交换方形 3.7.1，$c$ \Term{proper}{固有}，则对 $K\in\operatorname{ob}D^+(C)$，有\Term{canonical}{典范}\Term{arrow}{箭头}
\[f_*c_*c^!K\longrightarrow d_*d^!f_*K,\tag{6.5.3}\]
定义与 3.7.3 相同（\Term{isomorphism}{同构} $g_*c^!\xrightarrow\sim d^!f_*$ 在本情形也成立：\Term{arrow}{箭头}定义为\Term{base-change}{基变换}\Term{arrow}{箭头}的\Term{transpose}{转置}，再归结到 $d$ \Term{smooth}{光滑}和为\Term{closed-immersion}{闭浸入}的情形，验证其为\Term{isomorphism}{同构}）。取 $K=R\mathcal Hom_S(L_1,L_2)$，由 6.5.3、6.5.1 得到\Term{canonical}{典范}\Term{arrow}{箭头}
\SourcePage{137}
\[f_*c_*R\mathcal Hom(c_1^*L_1,c_2^!L_2)\longrightarrow d_*R\mathcal Hom(d_1^*(f_{1*}L_1),d_2^!(f_{2*}L_2)),\tag{6.5.4}\]
\[f_*:\operatorname{Hom}(c_1^*L_1,c_2^!L_2)\longrightarrow\operatorname{Hom}(d_1^*(f_{1*}L_1),d_2^!(f_{2*}L_2)),\tag{6.5.5}\]
分别类似于 3.7.5、3.7.6。
\textbf{6.6。} 对 $i=1,2$，设 $X_i$ 为\Term{finite-tor-dimension}{有限 Tor 维数}且 \Term{tor-independence}{Tor 无关}的 $S$ \Term{scheme}{概形}，$L_i\in\operatorname{ob}D(X_i)_{parf/S}$。如同 4.1，\Term{evaluation-map}{求值箭头} $DL_i\mathbin{\overset L\otimes}L_i\to K_{X_i}$ 的\Term{external-tensor-product}{外张量积}定义类似于 4.1.2、4.1.3、4.1.4 的“\Term{cup-product}{杯积}”\Term{pairing}{配对}。
设 $c:C\to X$、$d:D\to X$ 为\Term{proper}{固有}\Term{morphism}{态射}，$e=c\times_Xd:E\to X$，$P,Q\in\operatorname{ob}D(X)_{parf/S}$。由 6.1.1 有
\[c_*c^!P=R\mathcal Hom(Rc_*\mathcal O_C,P)\]
（对 $d$ 有类似\Term{isomorphism}{同构}）。设 $c,d$ 在 $e(E)$ 上方各点处\Term{finite-tor-dimension}{有限 Tor 维数}，则取\Term{tensor-product}{张量积}得到\Term{canonical}{典范}\Term{arrow}{箭头}
\[c_*c^!P\mathbin{\overset L\otimes}d_*d^!Q\longrightarrow R\mathcal Hom(Rc_*\mathcal O_C\mathbin{\overset L\otimes}Rd_*\mathcal O_D,P\mathbin{\overset L\otimes}Q).\tag{6.6.1}\]
当 $c,d$ \Term{tor-independence}{Tor 无关}时，由 Künneth 及 6.1.1，右端可写为 $e_*e^!(P\mathbin{\overset L\otimes}Q)$，于是 6.6.1 具有 4.2.3 的\Term{form}{形式}。取 $P=DL_1\mathbin{\overset L\otimes}_SL_2$、$Q=L_1\mathbin{\overset L\otimes}_SDL_2$，由 6.6.1 及上述\Term{pairing}{配对}，得到类似 4.2.4、4.2.5 的\Term{pairing}{配对}：
\[c_*R\mathcal Hom(c_1^*L_1,c_2^!L_2)\mathbin{\overset L\otimes}d_*R\mathcal Hom(d_2^*L_2,d_1^!L_1)\longrightarrow R\mathcal Hom(Rc_*\mathcal O_C\mathbin{\overset L\otimes}Rd_*\mathcal O_D,K_X),\tag{6.6.2}\]
\[\langle\ ,\ \rangle:\operatorname{Hom}(c_1^*L_1,c_2^!L_2)\otimes\operatorname{Hom}(d_2^*L_2,d_1^!L_1)\longrightarrow\operatorname{Hom}(Rc_*\mathcal O_C\mathbin{\overset L\otimes}Rd_*\mathcal O_D,K_X).\tag{6.6.3}\]
与 4.2 的\Term{pairing}{配对}一样，这些\Term{pairing}{配对}与\Term{etale-localization}{平展局部化}相容。当 $c,d$ \Term{tor-independence}{Tor 无关}时，6.6.2（相应地 6.6.3）的右端\Term{canonical}{典范}地等同于 $e_*K_E$（相应地 $H^0(E,K_E)$）。
\SourcePage{138}
\Term{theorem}{定理} 4.4 有如下类似结果。
\begin{SGAresult}{\Term{theorem}{定理}}{6.7}
对 $i=1,2$，设 $f_i:X_i\to Y_i$ 为\Term{proper}{固有} $S$ \Term{morphism}{态射}，$f=f_1\times_Sf_2:X\to Y$，$p_i:X\to X_i$、$q_i:Y\to Y_i$ 为\Term{projection}{投影}。再设有水平面为\Term{cartesian}{笛卡尔}方形的立方体 4.4.0，$c',c'',d',d''$ \Term{proper}{固有}，且 $c',c''$（相应地 $d',d''$）在 $c(C)$（相应地 $d(D)$）上方各点处\Term{finite-tor-dimension}{有限 Tor 维数}。置 $c_i'=p_ic'$、$c_i''=p_ic''$、$d_i'=q_id'$、$d_i''=p_id''$。最后设 $L_i\in\operatorname{ob}D(X_i)_{parf/S}$。设 $(X_1,X_2)$、$(X_1,Y_2)$、$(Y_1,X_2)$、$(Y_1,Y_2)$ 均在 $S$ 上 \Term{tor-independence}{Tor 无关}，$X_i,Y_i$ 在 $S$ 上有\Term{finite-tor-dimension}{有限 Tor 维数}。则有交换方形
\[\begin{tikzcd}[column sep=small]
f_*c'_*R\mathcal Hom(c_1'^*L_1,c_2'^!L_2)\mathbin{\overset L\otimes}f_*c''_*R\mathcal Hom(c_2''^*L_2,c_1''^!L_1)\arrow[r,"(1)"]\arrow[d,"(2)"']&f_*R\mathcal Hom(Rc'_*\mathcal O_{C'}\mathbin{\overset L\otimes}Rc''_*\mathcal O_{C''},K_X)\arrow[d,"(4)"]\\d'_*R\mathcal Hom(d_1'^*M_1,d_2'^!M_2)\mathbin{\overset L\otimes}d''_*R\mathcal Hom(d_2''^*M_2,d_1''^!M_1)\arrow[r,"(3)"]&R\mathcal Hom(Rd'_*\mathcal O_{D'}\mathbin{\overset L\otimes}Rd''_*\mathcal O_{D''},K_Y),
\end{tikzcd}\tag{6.7.1}\]
其中 $M_i=f_{i*}L_i$，(3) 由 6.6.2 给出；(1) 由对 6.6.2 应用 $f_*$ 并与\Term{canonical}{典范}\Term{arrow}{箭头} $f_*c'_*\mathbin{\overset L\otimes}f_*c''_*\to f_*(c'_*\mathbin{\overset L\otimes}c''_*)$ 复合得到；(2) 是与含 $f$ 的两个面对应的 6.5.4 \Term{arrow}{箭头}的\Term{tensor-product}{张量积}；(4) 是\Term{duality-isomorphism}{对偶同构} 6.1.1 与由\Term{canonical}{典范}\Term{arrow}{箭头}
\[Rd'_*\mathcal O_{D'}\mathbin{\overset L\otimes}Rd''_*\mathcal O_{D''}\longrightarrow f_*(Rc'_*\mathcal O_{C'}\mathbin{\overset L\otimes}Rc''_*\mathcal O_{C''})\]
定义的\Term{arrow}{箭头}的复合；该\Term{canonical}{典范}\Term{arrow}{箭头}\Term{adjunction}{伴随}于“\Term{base-change}{基变换}”\Term{arrow}{箭头} $f^*Rd'_*\mathcal O_{D'}\to Rc'_*\mathcal O_{C'}$、$f^*Rd''_*\mathcal O_{D''}\to Rc''_*\mathcal O_{C''}$ 的\Term{tensor-product}{张量积}。
\end{SGAresult}
证明实质上与 4.4 相同，少量所需修改留给读者。
当 $c',c''$ \Term{tor-independence}{Tor 无关}且 $d',d''$ 也 \Term{tor-independence}{Tor 无关}时，6.7.1 的\Term{arrow}{箭头} (4) 可写为 $f_*c_*K_C\to d_*K_D$，它仅由\Term{adjoint}{伴随箭头} $g_*K_C=g_*g^!K_D\to K_D$ 定义（与 4.4.1 的\Term{arrow}{箭头} (4) 一样）。不过，只考虑这一情形并不自然，因为这将排除超额相交的对应。
由 6.7 得到类似 4.5.1、4.7.1 的公式：在 6.7 的假设下，对 $u\in\operatorname{Hom}(c_1'^*L_1,c_2'^!L_2)$、$v\in\operatorname{Hom}(c_2''^*L_2,c_1''^!L_1)$，有
\SourcePage{139}
\[\langle f_*u,f_*v\rangle=f_*\langle u,v\rangle,\tag{6.7.2}\]
其中 $f_*u$（相应地 $f_*v$）为 6.5.5 意义下 $u$（相应地 $v$）的\Term{direct-image}{直像}对应，右端 $f_*$ 是\Term{homomorphism}{同态}
\[\operatorname{Hom}(Rc'_*\mathcal O_{C'}\mathbin{\overset L\otimes}Rc''_*\mathcal O_{C''},K_X)\longrightarrow\operatorname{Hom}(Rd'_*\mathcal O_{D'}\mathbin{\overset L\otimes}Rd''_*\mathcal O_{D''},K_Y),\tag{6.7.3}\]
它是\Term{duality-isomorphism}{对偶同构} 6.1.1 与由 6.7 所述\Term{canonical}{典范}\Term{arrow}{箭头} $Rd'_*\mathcal O_{D'}\mathbin{\overset L\otimes}Rd''_*\mathcal O_{D''}\to f_*(Rc'_*\mathcal O_{C'}\mathbin{\overset L\otimes}Rc''_*\mathcal O_{C''})$ 诱导的\Term{arrow}{箭头}的复合。当 $Y_1=Y_2=S=D'=D''$ 时，6.7.2 就是 Lefschetz--Verdier 公式。
作为说明，下面证明此公式蕴含 Woods Hole 公式 [1]。
\textbf{6.8。} 为此需要回顾：\Term{smooth}{光滑}\Term{scheme}{概形}的一个\Term{regular}{正则}\Term{embedding}{嵌入}的子\Term{scheme}{概形}所对应的“Hodge”\Term{cohomology-class}{上同调类}，以及 Grothendieck \Term{residue}{留数}（[8] III 9）。设 $f:X\to S$ 为\Term{smooth}{光滑}\Term{morphism}{态射}，纯\Term{relative-dimension}{相对维数}为 $N$，$i:Y\hookrightarrow X$ 为\Term{codimension}{余维} $d$ 的\Term{regular-immersion}{正则浸入}（即（SGA IV 16.9.2）、（SGA 6 VII 1）中，\Term{local-scheme}{局部}由 $d$ 个方程组成的\Term{regular-sequence}{正则序列}定义的闭子\Term{scheme}{概形}；用\Term{equivalence}{等价}术语（[8] III 1）、（SGA 6 VIII 1），是\Term{codimension}{余维} $d$ 的\Term{local-complete-intersection}{局部完全交}\Term{closed-immersion}{闭浸入}）。置 $g=fi$。如下给 $Y$ 配上一个元素
\[\operatorname{cl}_{X/S}(Y)\in\operatorname{Ext}_{\mathcal O_X}^d(\mathcal O_Y,\Omega^d_{X/S}),\tag{6.8.1}\]
称为与 $Y$ 对应的\Term{cohomology-class}{上同调类}。由“\Term{purity}{纯性}”\Term{theorem}{定理}（[8] III 7.2），有
\[\mathcal Ext_{\mathcal O_X}^i(\mathcal O_Y,\Omega^d_{X/S})=
\begin{cases}0&i\ne d,\\\mathcal Hom(\bigwedge^d\mathcal N_{Y/X},\Omega^d_{X/S}\otimes\mathcal O_Y)&i=d,\end{cases}\tag{6.8.2}\]
其中 $\mathcal N_{Y/X}=I/I^2$ 是余法\Term{sheaf}{层}，$I=\operatorname{Ker}(\mathcal O_X\to\mathcal O_Y)$。因此
\[\operatorname{Ext}_{\mathcal O_X}^d(\mathcal O_Y,\Omega^d_{X/S})=H^0(Y,\mathcal Hom(\bigwedge^d\mathcal N_{Y/X},\Omega^d_{X/S}\otimes\mathcal O_Y)).\tag{6.8.3}\]
\SourcePage{140}
\Term{differential}{微分} $d_{X/S}:\mathcal O_X\to\Omega^1_{X/S}$ 诱导 $\mathcal O_Y$ \Term{linear}{线性}映射 $d_{X/S}\otimes1:\mathcal N_{Y/X}\to\Omega^1_{X/S}\otimes\mathcal O_Y$。此映射的第 $d$ \Term{exterior-power}{外幂}是 $Y$ 上 $\bigwedge^d\mathcal N_{Y/X}^{\vee}\otimes\Omega^d_{X/S}$ 的全局\Term{section}{截面}；定义其经\Term{isomorphism}{同构} 6.8.3 在 $\operatorname{Ext}_{\mathcal O_X}^d(\mathcal O_Y,\Omega^d_{X/S})$ 中的像为类 6.8.1。
设 $g$ \Term{finite}{有限}且\Term{flat}{平坦}（所以 $d=N$）。取 $\omega\in H^0(Y,\bigwedge^N\mathcal N_{Y/X}^{\vee}\otimes\Omega^N_{X/S})$。由 6.8.2、6.1 可将 $\omega$ 看作\Term{homomorphism}{同态}
\[\mathcal O_Y\longrightarrow g^!\mathcal O_S=\bigwedge^N\mathcal N_{Y/X}^{\vee}\otimes\Omega^N_{X/S},\]
由\Term{adjunction}{伴随}性又可看作\Term{homomorphism}{同态} $g_*\mathcal O_Y\to\mathcal O_S$。定义 $\omega$（相对于 $g$）的\Term{residue}{留数}为 $1$ 经此\Term{homomorphism}{同态}在 $\mathcal O_S$ 中的像，记为 $\operatorname{Res}_{Y/S}(\omega)$（参见 [8] III 9）。另一方面，以 $i_*\omega\in H^N(X,\Omega^N_{X/S})$ 记 $\omega$ 经\Term{isomorphism}{同构} 6.8.3 与\Term{canonical}{典范}\Term{arrow}{箭头} $\operatorname{Ext}_{\mathcal O_X}^N(\mathcal O_Y,\Omega^N_{X/S})\to H^N(X,\Omega^N_{X/S})$ 的复合所得的像。当 $f$ \Term{proper}{固有}时，由\Term{trace-morphism}{迹态射}（即\Term{adjoint}{伴随箭头}）$i_*i^!\to1$、$f_*f^!\to1$ 的\Term{transitivity}{传递性}，有
\[\operatorname{Res}_{Y/S}(\omega)=\operatorname{Tr}_f(i_*\omega),\tag{6.8.4}\]
其中 $\operatorname{Tr}_f:H^N(X,\Omega^N_{X/S})\to H^0(S,\mathcal O_S)$ 由\Term{adjoint}{伴随箭头} $f_*f^!\mathcal O_S=Rf_*\Omega^N_{X/S}[N]\to\mathcal O_S$（6.1）定义。另一方面，无须\Term{properness}{固有性}假设，有重要的相容性（[8] III 9 R6）
\[\operatorname{Res}_{Y/S}(y\operatorname{cl}_{X/S}(Y))=\operatorname{Tr}_{Y/S}(y),\tag{6.8.5}\]
其中 $y\in H^0(Y,\mathcal O_Y)$，$\operatorname{Tr}_{Y/S}:g_*\mathcal O_Y\to\mathcal O_S$ 表示通常意义的迹，即 $\operatorname{Tr}_{Y/S}(y)$ 是 $g_*\mathcal O_Y$ 上乘以 $y$ 的\Term{endomorphism}{自同态}之迹（所引文献未证明此相容性；$N=1$ 的情形见 Raynaud，《\Term{henselian-local-ring}{Hensel 局部环}》，VIII 1，Lecture Notes in Math. 第 169 号，Springer Verlag，1970）。换言之，$\operatorname{Tr}_{Y/S}$ 通过\Term{adjunction}{伴随}性对应于由乘以 $\operatorname{cl}_{X/S}(Y)$ 定义的\Term{homomorphism}{同态} $\mathcal O_Y\to g^!\mathcal O_S$。所以对 $F\in\operatorname{ob}D(S)_{coh}$，由乘以 $\operatorname{cl}_{X/S}(Y)$ 给出的\Term{arrow}{箭头}
\[g^*F\longrightarrow g^!F=g^*F\mathbin{\overset L\otimes}g^!\mathcal O_S\tag{6.8.6}\]
通过\Term{adjunction}{伴随}性对应于\Term{arrow}{箭头}
\SourcePage{141}
\[g_*g^*F=F\mathbin{\overset L\otimes}g_*\mathcal O_Y\longrightarrow F,\tag{6.8.7}\]
后者由 $\operatorname{Tr}_{Y/S}:g_*\mathcal O_Y\to\mathcal O_S$ 定义。
\textbf{6.9。} 设 $X$（相应地 $Y$）为\Term{relative-dimension}{相对维数} $m$（相应地 $n$）的\Term{proper}{固有}\Term{smooth}{光滑} $S$ \Term{scheme}{概形}，$Z=X\times_SY$，$a:A\hookrightarrow Z$、$b:B\hookrightarrow Z$ 为\Term{regular-immersion}{正则浸入}，使得 $a_2:A\to Y$、$b_1:B\to X$ \Term{finite}{有限}且\Term{flat}{平坦}（以 $p_1:Z\to X$、$p_2:Z\to Y$ 记\Term{projection}{投影}，$a_i=p_ia$、$b_i=p_ib$）。设 $c:C=A\times_ZB\to Z$，$L\in\operatorname{ob}D(X)_{parf}$、$M\in\operatorname{ob}D(Y)_{parf}$（下标 $parf$ 表示“绝对意义下\Term{perfect}{完美}”；由于 $X,Y$ \Term{smooth}{光滑}，这也\Term{equivalence}{等价}于相对于 $S$ \Term{perfect}{完美}）。设 $C$ 在 $S$ 上\Term{finite}{有限}且\Term{flat}{平坦}。这些假设蕴含 $c$ 为\Term{regular-immersion}{正则浸入}，且 $a,b$ \Term{tor-independence}{Tor 无关}。
有 $p_1^*\Omega^1_{X/S}=\Omega^1_{Z/Y}$、$p_2^*\Omega^1_{Y/S}=\Omega^1_{Z/X}$，及\Term{resolution}{分解} $\Omega^1_{X/S}=\Omega^1_{Z/Y}\oplus\Omega^1_{Z/X}$，从而对任意整数 $r$，有\Term{resolution}{分解}
\[\Omega^r_{Z/S}=\bigoplus_{i+j=r}\Omega^i_{Z/Y}\otimes\Omega^j_{Z/X}.\tag{6.9.0}\]
另一方面，由 6.1.2，有 $p_2^!M=p_2^*M\otimes\Omega^m_{Z/Y}[m]$，故
\[H^0(Z,a_*a^!R\mathcal Hom(p_1^*L,p_2^!M))=\operatorname{Ext}_{\mathcal O_Z}^m(a_1^*L,a_2^*M\otimes\Omega^m_{Z/Y}),\]
由 6.4.4、6.8.2，得到\Term{isomorphism}{同构}
\[\operatorname{Hom}(a_1^*L,a_2^!M)=H^0(A,R\mathcal Hom(a_1^*L,a_2^*M)\otimes\bigwedge^m\mathcal N_{A/Z}^{\vee}\otimes\Omega^m_{Z/Y}).\tag{6.9.1}\]
同样有\Term{isomorphism}{同构}
\[\operatorname{Hom}(b_2^*M,b_1^!L)=H^0(B,R\mathcal Hom(b_2^*M,b_1^*L)\otimes\bigwedge^n\mathcal N_{B/Z}^{\vee}\otimes\Omega^n_{Z/X}).\tag{6.9.2}\]
取 $u\in\operatorname{Hom}(a_1^*L,a_2^*M)$、$v\in\operatorname{Hom}(b_2^*M,b_1^*L)$，以
\[\langle u,v\rangle\in H^0(C,\mathcal O_C)\tag{6.9.3}\]
\SourcePage{142}
记\Term{section}{截面} $\operatorname{Tr}(u_Cv_C)=\operatorname{Tr}(v_Cu_C)$，其中 $u_C\in\operatorname{Hom}(c_1^*L,c_2^*M)$、$v_C\in\operatorname{Hom}(c_2^*M,c_1^*L)$ 分别由 $u,v$ 诱导。考虑由 $u$（相应地 $v$）与迹\Term{homomorphism}{同态} 6.8.7 $a_{2*}a_2^*M\to M$（相应地 $b_{1*}b_1^*L\to L$）复合所得的\Term{cohomological-correspondence}{上同调对应} $u^!\in\operatorname{Hom}(a_1^*L,a_2^!M)$（相应地 $v^!\in\operatorname{Hom}(b_2^*M,b_1^!L)$）；\Term{equivalence}{等价}地（6.8.6），它由乘以 $\operatorname{cl}_{Z/Y}(A)$（相应地 $\operatorname{cl}_{Z/X}(B)$）得到。顺便指出，$\operatorname{cl}_{Z/Y}(A)$（相应地 $\operatorname{cl}_{Z/X}(B)$）是 $\operatorname{cl}_{Z/S}(A)$（相应地 $\operatorname{cl}_{Z/S}(B)$）经 6.9.0 的\Term{projection}{投影} $\Omega^m_{Z/S}\to\Omega^m_{Z/Y}$（相应地 $\Omega^n_{Z/S}\to\Omega^n_{Z/X}$）所定义的 $H^0(A,\bigwedge^m\mathcal N_{A/Z}^{\vee}\otimes\Omega^m_{Z/Y})$（相应地 $H^0(B,\bigwedge^n\mathcal N_{B/Z}^{\vee}\otimes\Omega^n_{Z/X})$）中的\Term{set-component}{分量}。容易验证，6.6.3 定义的\Term{cup-product}{杯积}
\[\langle u^!,v^!\rangle\in\operatorname{Ext}_{\mathcal O_Z}^{m+n}(\mathcal O_C,\Omega^{m+n}_{Z/S})=H^0(C,\bigwedge^{m+n}\mathcal N_{C/Z}^{\vee}\otimes\Omega^{m+n}_{Z/S})\]
由下式给出：
\[\langle u^!,v^!\rangle=\langle u,v\rangle\operatorname{cl}_{Z/Y}(A)\operatorname{cl}_{Z/X}(B).\tag{6.9.4}\]
另一方面可验证，\Term{homomorphism}{同态} 6.7.3
\[\operatorname{Ext}_{\mathcal O_Z}^{m+n}(\mathcal O_C,\Omega^{m+n}_{Z/S})\longrightarrow H^0(S,\mathcal O_S)\]
是\Term{canonical}{典范}\Term{arrow}{箭头} $\operatorname{Ext}_{\mathcal O_Z}^{m+n}(\mathcal O_C,\Omega^{m+n}_{Z/S})\to H^{m+n}(Z,\Omega^{m+n}_{Z/S})$ 与\Term{trace-morphism}{迹态射} $\operatorname{Tr}_{Z/S}:H^{m+n}(Z,\Omega^{m+n}_{Z/S})\to H^0(S,\mathcal O_S)$ 的复合。由 Lefschetz--Verdier 公式 6.7.2 得到关系
\[\langle(u^!)_*,(v^!)_*\rangle=\operatorname{Tr}_{Z/S}(\langle u,v\rangle\operatorname{cl}_{Z/Y}(A)\operatorname{cl}_{Z/X}(B)).\tag{6.9.5}\]
借助 6.8.4，可把右端写成\Term{residue}{留数}；综上得到：
\begin{SGAresult}{\Term{theorem}{定理}}{6.10}
在 6.9 第一段的假设下，设 $u\in\operatorname{Hom}(a_1^*L,a_2^*M)$、$v\in\operatorname{Hom}(b_2^*M,b_1^*L)$。考虑\Term{support}{支撑}在 $A$（相应地 $B$）中的从 $L$ 到 $M$（相应地 $M$ 到 $L$）的\Term{cohomological-correspondence}{上同调对应} $u^!$（相应地 $v^!$），它由 $u$（相应地 $v$）与迹 $a_{2*}a_2^*M\to M$（相应地 $b_{1*}b_1^*L\to L$）（6.8.7）复合得到；考虑取 $u^!$（相应地 $v^!$）的\Term{direct-image}{直像}（6.5.5）所得的\Term{homomorphism}{同态} $(u^!)_*:f_*L\to g_*M$（相应地 $(v^!)_*:g_*M\to f_*L$），其中 $f:X\to S$、$g:Y\to S$ 为\Term{projection}{投影}。
\SourcePage{143}
则用 6.8、6.9.3 的记号，有
\[\langle(u^!)_*,(v^!)_*\rangle=\operatorname{Res}_{C/S}(\langle u,v\rangle\operatorname{cl}_{Z/Y}(A)\operatorname{cl}_{Z/X}(B)).\tag{6.10.1}\]
\end{SGAresult}
\textbf{注记 6.11。} a) 如同 3.6，利用 3.5 b) 对对应的\Term{direct-image}{直像}的描述，可验证：若 $h:Z\to S$ 为\Term{projection}{投影}，则 $(u^!)_*$（相应地 $(v^!)_*$）为下列\Term{arrow}{箭头}的复合：
\[f_*L\xrightarrow{(1)}(ha)_*a_1^*L\xrightarrow{(ha)_*u}(ha)_*a_2^*M=g_*a_{2*}a_2^*M\xrightarrow{(2)}g_*M\]
（相应地 $g_*M\xrightarrow{(1)}(hb)_*b_2^*M\xrightarrow{(hb)_*v}(hb)_*b_1^*L=f_*b_{1*}b_1^*L\xrightarrow{(2)}f_*L$），其中 (1) 是通常的\Term{functoriality}{函子性}\Term{arrow}{箭头}，由\Term{adjoint}{伴随箭头} $1\to a_{1*}a_1^*$（相应地 $1\to b_{2*}b_2^*$）定义；(2) 由迹 $a_{2*}a_2^*M\to M$（相应地 $b_{1*}b_1^*L\to L$）（6.8.7）给出。

b) 设 $S$ 为域的谱，则交 $C=A\times_{X\times_SY}B$ 由\Term{isolated-point}{孤立点}组成。在 $C$ 的每个点 $z$ 的像附近，选择 $A,B$ 的方程的\Term{regular-sequence}{正则序列} $(s_1,\ldots,s_m)$、$(t_1,\ldots,t_n)$。用（[8] III 9）的\Term{residue}{留数}\Term{signature}{符号}记号，6.10.1 写成
\[\langle(u^!)_*,(v^!)_*\rangle=\sum_{z\in C}\operatorname{Res}_z\left(\langle u,v\rangle\prod_{1\leq i\leq m}(d_{Z/Y}s_i/s_i)\prod_{1\leq j\leq n}(d_{Z/X}t_j/t_j)\right),\tag{6.11.1}\]
其中 $d_{Z/Y}s_i$（相应地 $d_{Z/X}t_j$）表示 $d_{Z/S}s_i$（相应地 $d_{Z/S}t_j$）在\Term{resolution}{分解} 6.9.0 中的 $(1,0)$ 型（相应地 $(0,1)$ 型）\Term{set-component}{分量}。
\begin{SGAresult}{推论}{6.12（Woods Hole 公式 [1]）}
设 $S$ 为域的谱，$X$ 为\Term{proper}{固有}\Term{smooth}{光滑} $S$ \Term{scheme}{概形}，$F$ 为 $X$ 的 $S$ \Term{endomorphism}{自同态}，$L\in\operatorname{ob}D(X)_{parf}$（即 $X$ 上的\Term{perfect-complex}{完美复形}），$u\in\operatorname{Hom}(F^*L,L)$。设\Term{fixed-point}{不动点}\Term{scheme}{概形} $X^F$ \Term{finite}{有限}且由\Term{transverse}{横截}点组成（即 $F$ 的图与\Term{diagonal}{对角线}在这些点\Term{transverse-intersection}{横截相交}）。则
\[\sum_i(-1)^i\operatorname{Tr}((F,u),H^i(X,L))=\sum_{x\in X^F}\operatorname{Tr}(u_x)/\det(1-dF_x),\]
其中 $(F,u)$ 是 $H^i(X,L)\to H^i(X,F^*L)$
\SourcePage{144}
（\Term{inverse-image}{逆像}）与 $H^i(X,u)$ 复合所得的 $H^i(X,L)$ 的\Term{endomorphism}{自同态}；$\operatorname{Tr}(u_x)$ 表示 $u$ 在 $L$ 的 $x$ 处\Term{fiber}{纤维}上诱导的\Term{endomorphism}{自同态}的迹在 $x$ 处的值；$\det(1-dF_x)$ 表示映射“$1$ 减去 $F$ 在 $\Omega^1_{X,x}$ 上诱导的映射”的\Term{determinant}{行列式}在 $x$ 处的值。
\end{SGAresult}
\begin{proof}
结合 6.11 a)、b) 应用 6.10。置 $Z=X\times_SX$，取 $x\in X^F$。在 $x$ 附近选择\Term{local-coordinate}{局部坐标}系 $(x_1,\ldots,x_n)$（定义到 $\mathbb A_S^n$ 的\Term{etale}{平展}\Term{morphism}{态射}），置 $F_i=x_iF$；以 $(x_1,\ldots,x_n,y_1,\ldots,y_n)$ 记 $Z$ 在 $z=(x,x)$ 附近的相应\Term{local-coordinate}{局部坐标}，$x_i=x_i\otimes1$、$y_i=1\otimes x_i$。在 $z$ 附近，$(x_1-y_1,\ldots,x_n-y_n)$（相应地 $(y_1-F_1,\ldots,y_n-F_n)$）是\Term{diagonal}{对角线} $A$（相应地 $F$ 的图 $B$）的一组\Term{regular}{正则}方程。6.11.1 右端在 $z$ 处的贡献为
\[\langle u,1\rangle_z=\operatorname{Res}_z\left(\operatorname{Tr}(u)\frac{dx_1\cdots dx_n\,dy_1\cdots dy_n}{(x_1-y_1)\cdots(x_n-y_n)(y_1-F_1)\cdots(y_n-F_n)}\right).\]
注意分子中可以用 $dy_i-dF_i=dy_i-\sum_{1\leq j\leq n}(\partial F_i/\partial x_j)dx_j$ 代替 $dy_i$。将（[8] III 9 R3）应用于情形 $B\hookrightarrow Z\to S$、$(s_1,\ldots,s_p)=(y_1-F_1,\ldots,y_n-F_n)$、$(t_1,\ldots,t_n)=(x_1-y_1,\ldots,x_n-y_n)$、$\omega=\operatorname{Tr}(u)dx_1\cdots dx_n$，得到
\[\langle u,1\rangle_z=\operatorname{Res}_x\left(\operatorname{Tr}(u)\frac{dx_1\cdots dx_n}{(x_1-F_1)\cdots(x_n-F_n)}\right).\tag{*}\]
由于 $x$ 为\Term{transverse}{横截}点，$x_i-F_i$ 构成 $\mathcal O_{X,x}$ 的\Term{regular-system-of-parameters}{正则参数系}，故在 $x$ 附近可写为
\[x_i=\sum_j(x_j-F_j)c_{ij},\]
其中 $(c_{ij})$ 是由 $\mathcal O_X$ 的\Term{section}{截面}组成、在 $x$ 处\Term{invertible}{可逆}的矩阵，其\Term{determinant}{行列式}在 $x$ 处的值为 $\det(1-dF_x)^{-1}$。由（[8] III 9 R1）及 (*) 得
\[\langle u,1\rangle_z=\operatorname{Res}_x\left(\operatorname{Tr}(u)\det(c_{ij})\frac{dx_1\cdots dx_n}{x_1\cdots x_n}\right).\]
再由（[8] III 9 R6），
\SourcePage{145}
\[\langle u,1\rangle_z=\operatorname{Tr}(u(x))\det(c_{ij}(x))=\operatorname{Tr}(u(x))\det(1-dF_x)^{-1}.\]
故推论由 6.11.1 得证。
\end{proof}
\textbf{注记 6.12.1。} 若仅假设 $X^F$ \Term{finite}{有限}，上面的计算表明
\[\sum_i(-1)^i\operatorname{Tr}((F,u),H^i(X,L))=\sum_{x\in X^F}\operatorname{Res}_x\left(\operatorname{Tr}(u)\frac{dx_1\cdots dx_n}{(x_1-F_1)\cdots(x_n-F_n)}\right),\]
其中 $(x_i)$ 是 $x$ 处的一组\Term{local-coordinate}{局部坐标}，$F_i$ 为 $F$ 在此坐标系中（$x$ 附近）的坐标。另一方面，读者可以确信：若 $X$ \Term{proper}{固有}且仅假设在 $X^F$ 的各点\Term{smooth}{光滑}，上式仍成立。
\textbf{6.13。} 设 $S$ 为\Term{finite-field}{有限域} $\mathbb F_q$ 的一个\Term{algebraic-closure}{代数闭包}的谱，$X/S$ 由\Term{proper}{固有}\Term{smooth}{光滑}\Term{scheme}{概形} $X_0/\mathbb F_q$ 经\Term{scalar-extension}{标量扩张}而来。取 $F$ 为 $X/S$ 上“由坐标取 $q$ 次幂定义”的 Frobenius \Term{endomorphism}{自同态}。有 $X^F=X_0(\mathbb F_q)$，且由于 $dF=0$，$X^F$ 由\Term{transverse}{横截}点组成。若 $L$ 是 $X$ 上\Term{locally-free}{局部自由}\Term{finite-type}{有限型}\Term{sheaf}{层}，$u\in\operatorname{Hom}(F^*L,L)$，则 Woods Hole 公式写为
\[\sum_i(-1)^i\operatorname{Tr}((F,u),H^i(X,L))=\sum_{x\in X_0(\mathbb F_q)}\operatorname{Tr}(u_x,L_x).\tag{6.13.1}\]
特别地，取 $L=\mathcal O_X$、$u$ 为\Term{canonical}{典范}\Term{isomorphism}{同构} $F^*\mathcal O_X\xrightarrow\sim\mathcal O_X$，得到
\[\sum_i(-1)^i\operatorname{Tr}(F,H^i(X,\mathcal O_X))=\operatorname{Card}(X_0(\mathbb F_q))\pmod p.\tag{6.13.2}\]
Artin--Schreier 理论可把左端写成 $\sum_i(-1)^i\operatorname{Tr}(F,H^i(X,\mathbb F_p))$，其中 $H^*(X,\mathbb F_p)$ 表示 $X$ 的取值于常值\Term{sheaf}{层} $\mathbb F_p$ 的\Term{etale-cohomology}{平展上同调}。精细化此论证，Deligne 在（SGA $4\frac12$ Fonctions $L$）中由 6.13.1 得到 6.13.2 的以下推广：
\SourcePage{146}
\begin{SGAresult}{\Term{theorem}{定理}}{6.13.3（Deligne）}
设 $X_0$ 是 $\mathbb F_q$ 上\Term{separated}{分离}\Term{finite-type}{有限型}\Term{scheme}{概形}，$X$ 为 $X_0$ 经\Term{scalar-extension}{标量扩张}到 $\mathbb F_q$ 的\Term{algebraic-closure}{代数闭包} $k$ 所得的\Term{scheme}{概形}，$F$ 为 $X/k$ 的 Frobenius \Term{endomorphism}{自同态}。则对任意 $L_0\in\operatorname{ob}D_{ctf}(X_0,\mathbb F_p)$，以 $L$ 记 $L_0$ 在 $X$ 上的\Term{inverse-image}{逆像}，有
\[\sum_i(-1)^i\operatorname{Tr}(F,H_c^i(X,L))=\sum_{x\in X^F}\operatorname{Tr}(F,L_x).\]
\end{SGAresult}
指出，6.13.3 中 $L_0=\mathbb F_p$ 的特例由 Katz 在（SGA 7 XXII）中用完全不同的方法建立，其基础是对\Term{hypersurface}{超曲面}的 zeta 函数作 Dwork 式计算。另一方面，在\Term{proper}{固有}\Term{smooth}{光滑}情形，6.13.2 也是\Term{crystalline-cohomology}{晶体上同调}的 Lefschetz 公式（[5] VII 3.1.9）的直接推论。事实上，\Term{crystalline-cohomology}{晶体上同调}给出更精确的信息：zeta 函数的每个因子 $\det(1-Ft,H^i(X_0/W))\in W[t]$ 模 $p$ 同余于 $\det(1-Ft,H^i(X_0,\mathcal O_{X_0}))$（至少在 $H^i(X_0/W)$ \Term{torsion-free}{无挠}时）。Katz 在 [11] 中表明，对直至某一秩的 Hodge 数（在所关心的\Term{dimension}{维数}）为零的\Term{smooth}{光滑}\Term{hypersurface}{超曲面}（或\Term{complete-intersection}{完全交}），可借助 B. Mazur [14] 的估计，得到更高阶的同余（涉及 $H^i(X,\Omega_X^j)$ 及 Cartier \Term{operator}{算子}）。
另一方面，若 $X_0$ 是 $\mathbb F_q$ 上\Term{proper}{固有}且\Term{geometrically-integral}{几何整}的\Term{scheme}{概形}，并且 $H^i(X_0,\mathcal O_{X_0})=0$（$i>0$），则将 Lefschetz--Verdier 公式 6.7.2 应用于 $X=X_0\otimes k$ 和 $X$ 的 Frobenius \Term{endomorphism}{自同态}，便得 $X_0$ 至少有一个 $\mathbb F_q$ \Term{rational-point}{有理点}（因为 $\operatorname{Tr}(F,R\Gamma(X,\mathcal O_X))=1$）。Serre 问：若将 $\mathbb F_q$ 换成满足 $(C_1)$ 的域（[17] II 3），例如 $\mathbb C(t)$（Tsen），或 $\mathbb C((t))$（Lang）（所引文献还有其他例子），这一结论是否仍成立。他还指出：若 $X_0$ 是 $\mathbb P^n$ 中\Term{multidegree}{多重次数}为 $(m_1,\ldots,m_r)$ 的\Term{complete-intersection}{完全交}，则\Term{cohomologically-trivial}{上同调平凡}的条件\Term{equivalence}{等价}于（[16] 第 78 号）$m_1+\cdots+m_r<n+1$，而 $(C_1)$ 条件的陈述中出现的正是这个条件。
\SourcePage{147}
\textbf{6.14。} Woods Hole 公式还有许多其他应用，参见 [2] 与 Beauville 的报告 [4]。
\textbf{6.15。} 撰稿人没有研究\Term{dimension}{维数} $\geq1$ 的\Term{fixed-point}{不动点}分支的情形；但或许可以由 6.7.2（可能还须结合 Riemann--Roch）\Term{derive}{导出} Donovan [7]、Iversen [10]、Nielsen [15] 风格的 Lefschetz--\Term{riemann-roch-formula}{Riemann--Roch 公式}（至少其值取在\Term{base-field}{基域}中）。另一方面，将 6.7.2 应用于\Term{duality}{对偶}数上，也可能给出类似 Bott [6]、[3] 的\Term{vector-field}{向量场}\Term{residue}{留数}公式。
\SourcePage{148}
\SGAsection{}{参考文献}
\begin{description}
\item[{[1]}] Atiyah, M. 与 Bott, R.，《关于 Woods Hole \Term{fixed-point}{不动点}\Term{theorem}{定理}》（On the Woods Hole Fixed point theorem），Proceedings of the Woods Hole Conference on Algebraic Geometry，1964。
\item[{[2]}] Atiyah, M. 与 Bott, R.，《\Term{elliptic-complex}{椭圆复形}的 Lefschetz \Term{fixed-point}{不动点}\Term{theorem}{定理}札记》（Notes on the Lefschetz fixed point theorem for elliptic complexes），Harvard University，1964。
\item[{[3]}] Baum, P. 与 Bott, R.，《\Term{meromorphic}{亚纯}\Term{vector-field}{向量场}的零点》（On the zeroes of meromorphic vector-fields），收入 Essays on topology and related topics，献给 Georges de Rham 的纪念文集，Springer Verlag，1970。
\item[{[4]}] Beauville, A.，《\Term{coherent}{凝聚}\Term{cohomology}{上同调}的\Term{fixed-point}{不动点}公式》（Formules de points fixes en cohomologie cohérente），收入 Séminaire de géométrie algébrique，Université de Paris-Sud，Orsay，1970。
\item[{[5]}] Berthelot, P.，《\Term{characteristic}{特征} $p>0$ 的\Term{scheme}{概形}的\Term{crystalline-cohomology}{晶体上同调}》（Cohomologie cristalline des schémas de car. $p>0$），Lecture Notes in Mathematics，第 407 号，Springer-Verlag，1974。
\item[{[6]}] Bott, R.，《\Term{holomorphic}{全纯}\Term{vector-field}{向量场}的\Term{residue}{留数}公式》（A residue formula for holomorphic vector-fields），J. of Diff. Geometry，第 1 卷，第 311--330 页，1967。
\item[{[7]}] Donovan, P.，《Lefschetz--\Term{riemann-roch-formula}{Riemann--Roch 公式}》（The Lefschetz-Riemann-Roch formula），Bull. S.M.F.，97，第 257--273 页，1969。
\item[{[8]}] Hartshorne, R.，《\Term{residue}{留数}与\Term{duality}{对偶}》（Residues and duality），Lecture Notes in Mathematics，第 20 号，Springer-Verlag，1966。
\item[{[9]}] Illusie, L.，《\Term{cotangent-complex}{余切复形}与\Term{deformation}{形变} I》（Complexe cotangent et déformations I），Lecture Notes in Math.，第 239 号，Springer-Verlag，1971。
\item[{[10]}] Iversen, B.，《\Term{torus}{环面}\Term{action}{作用}于\Term{algebraic-variety}{代数簇}的\Term{fixed-point}{不动点}公式》（A fixed point formula for action of tori on algebraic varieties），Inv. Math.，16，第 229--236 页，1972。
\item[{[11]}] Katz, N.，《致 B. Mazur 的信》（Lettre à B. Mazur），1971。
\SourcePage{149}
\item[{[12]}] Kleiman, S.，《\Term{algebraic-cycle}{代数圈}与 Weil 猜想》（Algebraic cycles and the Weil conjectures），收入 Dix exposés sur la cohomologie des schémas，North Holland Pub. Company，1968。
\item[{[13]}] Langlands, R. P.，《模\Term{form}{形式}与 $\ell$ 进表示》（Modular forms and $\ell$-adic representations），收入 Modular functions of one variable II，Lecture Notes in Math.，第 349 号，Springer-Verlag，1973。
\item[{[14]}] Mazur, B.，《Frobenius 与 Hodge \Term{filtered}{滤过}（估计）》（Frobenius and the Hodge filtration (estimates)），Ann. of Math.，98，第 58--95 页，1973。
\item[{[15]}] Nielsen, H. A.，《可\Term{diagonal}{对角线}性化的\Term{coherent-sheaf}{凝聚层}》（Diagonalizably linearized coherent sheaves），Bull. SMF，102，第 85--97 页，1974。
\item[{[16]}] Serre, J.-P.，《\Term{coherent-algebraic-sheaf}{凝聚代数层}》（Faisceaux algébriques cohérents），Ann. of Math.，61，第 197--278 页，1955。
\item[{[17]}] Serre, J.-P.，《Galois \Term{cohomology}{上同调}》（Cohomologie galoisienne），Lecture Notes in Math.，第 5 号，Springer Verlag，1964。
\item[{[18]}] Verdier, J.-L.，《\Term{etale-cohomology}{平展上同调}的 Lefschetz \Term{fixed-point}{不动点}公式》（The Lefschetz fixed point formula in étale cohomology），Proceedings of a conference on local fields，Springer-Verlag，1967。
\end{description}
SGA $4\frac12$：《\Term{etale-cohomology}{平展上同调}》（Cohomologie étale），P. Deligne 著，Séminaire de Géométrie Algébrique du Bois-Marie，Lecture Notes in Mathematics，第 569 号，Springer-Verlag，1977。
其中：
\begin{description}
\item[Finitude：] 《$\ell$ 进\Term{cohomology}{上同调}的\Term{finiteness-theorem}{有限性定理}》（Théorème de finitude en cohomologie $\ell$-adique）。
\item[Cycle：] 《与一个圈对应的\Term{cohomology-class}{上同调类}》（La classe de cohomologie associée à un cycle），A. Grothendieck 著，P. Deligne 撰写（代替缺失的 SGA 5 第 IV 讲）。
\item[Rapport：] 《迹公式报告》（Rapport sur la formule des traces）。
\item[Fonctions $L$：] 《模 $\ell^n$ 的 $L$ 函数与模 $p$ 的 $L$ 函数》（Fonctions $L$ mod. $\ell^n$ et mod. $p$）。
\end{description}
